% Authoritative module; see specific_highlights_index.md.
\ifdefined\HighlightsTenTermTwoLoaded\endinput\fi
\let\HighlightsTenTermTwoLoaded\relax
\HighlightDeclareConcept{outcome-space}{
  role={outcome},
  notation={$\GTenHLOutcome{\omega}$, $\GTenHLOutcome{\Omega}$},
  name={\GTenHLOutcome{Outcome} \GTenHLOutcome{and} \GTenHLOutcome{sample} \GTenHLOutcome{space}},
  brief={One complete result; all possible complete results.},
  definition={$\omega$ is one complete outcome; $\Omega$ is the set of all possible outcomes.},
  source={book/grade10/term_2/1-practice_activities}
}

\HighlightDeclareConcept{event}{
  role={event},
  notation={$\HLMathEvent{A}$, $\HLMathEvent{E}$, $\HLMathEvent{E^c}$; $\{\GTenHLValue{X}\leq \GTenHLValue{a}\}$},
  name={\HLMathEvent{Event} \HLMathEvent{and} \HLMathEvent{complement}},
  brief={Outcomes meeting a condition; those outside it.},
  definition={$A$ and $E$ are sets of outcomes; $E^c$ contains those outside $E$. A condition such as $X\leq a$ also defines an event; $B$ is another event name used in the activities.},
  source={book/grade10/term_2/1-practice_activities}
}

\HighlightDeclareConcept{discrete-support}{
  role={quantity.primary},
  notation={$\GTenHLValue{S_X}=\{\GTenHLValue{x}:\HLProbability{P}(\GTenHLValue{X}=\GTenHLValue{x})>0\}$},
  name={\GTenHLValue{Support}},
  brief={Distinct values with positive probability, in the variable's units.},
  definition={$S_X$ contains the distinct values with positive mass; $S_T$ is the support of $T$. Only values with positive probability belong to the discrete support, including at parameter endpoints.},
  source={book/grade10/term_2/1-practice_activities}
}

\HighlightDeclareConcept{second}{
  role={quantity.secondary},
  notation={$\GTenHLSecond{Y}$, $\GTenHLSecond{y}$; $\GTenHLSecond{Y}=f(\GTenHLValue{X})$; $\GTenHLSecond{Y_i}=f_i(\GTenHLValue{X})$},
  name={\GTenHLSecond{Second} \GTenHLSecond{or} \GTenHLSecond{transformed} \GTenHLSecond{variable}},
  brief={Another numerical description, with its own units.},
  definition={$Y$, $y$ describe another quantity and its values; $f$ or $f_i$ gives the transformation rule. $E$, $C$ or $L$ can also name a second variable: its meaning and units depend on the problem.},
  source={book/grade10/term_2/1-practice_activities}
}

\HighlightDeclareConcept{third}{
  role={quantity.tertiary},
  notation={$\GTenHLThird{Z}$, $\GTenHLThird{z}$, $\GTenHLThird{C}$},
  name={\GTenHLThird{Third} \GTenHLThird{variable}},
  brief={A third numerical description; an indicator is unitless.},
  definition={$Z$ or $C$ represents a third quantity and $z$ a possible value. An indicator equals 1 when its condition holds and 0 otherwise; it has no units. Here $C$ is not a confidence level.},
  source={book/grade10/term_2/1-practice_activities}
}

\HighlightDeclareConcept{mathematical-operations}{
  role={context},
  notation={$\sum$, $\binom nk$, $n!$, $e$, $\sim$},
  name={\HLMathContext{Operations} \HLMathContext{and} \HLMathContext{distribution} \HLMathContext{notation}},
  brief={Sum, count selections, take factorials, use the exponential constant, or name a model.},
  definition={$\sum$ adds terms; $\sim$ means ``is distributed as''. Grade 10 also uses $\binom nk$ for selections, factorials ($0!=1$), and exponential base $e$. Grade 11 uses $\sum x_i$ and $\sum x_i^2$ to summarise data.},
  source={book/grade10/term_2/1-practice_activities}
}

\HighlightDeclareConcept{assume}{
  role={context},
  notation={Independence, constant $p$, fixed observation rule},
  name={\HLMathContext{Model} \HLMathContext{conditions}},
  brief={Conditions to check before using the selected distribution.},
  definition={Independence and a common probability $p$ underpin repeated Bernoulli trials. Waiting-time models specify a stopping rule. Poisson requires a constant rate, independent counts in disjoint intervals and isolated events.},
  source={book/grade10/term_2/1-practice_activities}
}

\HighlightDeclareConcept{cutoff}{
  role={quantity.primary},
  notation={$\GTenHLValue{a}$, $\GTenHLValue{b}$, $\GTenHLValue{m}$, $\GTenHLValue{\lfloor x\rfloor}$},
  name={\GTenHLValue{Cutoff} \GTenHLValue{and} \GTenHLValue{integer} \GTenHLValue{part}},
  brief={A boundary in the variable's units; specify included endpoints.},
  definition={$a,b,m$ or $k$ set bounds or comparison values in the variable's units. $\lfloor x\rfloor$ is the greatest integer less than or equal to $x$. State whether each endpoint is included.},
  source={book/grade10/term_2/1-practice_activities}
}

\HighlightDeclareConcept{aux}{
  role={context},
  notation={$K,a,b,c$; $a_m,b_m,c_m$; $N_k,b_j,R_m$},
  name={\HLMathContext{Auxiliary} \HLMathContext{counts} \HLMathContext{and} \HLMathContext{states}},
  brief={Problem-defined counts, recurrence states, balances or survival quantities.},
  definition={$K,a,b,c$, their versions $a_m,b_m,c_m$, and $N_k,b_j,R_m$ denote counts, recurrence states, balances or survival quantities defined in each problem. Subscripts identify a stage or category, rather than an operation.},
  source={book/grade10/term_2/1-practice_activities}
}

\HighlightDeclareConcept{break}{
  role={probability.success},
  notation={$\HLSuccessProbability{p_i^*}$},
  name={\HLSuccessProbability{Individual} \HLSuccessProbability{break-even} \HLSuccessProbability{probability}},
  brief={Probability at which the specified expected net payoff is zero.},
  definition={$p_i^*$ is the success probability that makes the expected net profit of alternative $i$ zero. It is a break-even threshold, rather than a normal critical value.},
  source={book/grade10/term_2/1-practice_activities}
}

\HighlightDeclareConcept{transform}{
  role={context},
  notation={$a,b$; $f_i(0),f_i(1)$},
  name={\HLMathContext{Transformation} \HLMathContext{inputs}},
  brief={Scale, fixed addition or state payoffs, with required units.},
  definition={In $Y=aX+b$, $a$ is the scale factor and $b$ the fixed term. $f_i(0),f_i(1)$ are the payoffs or values assigned to the two states. Units must be compatible with the resulting variable.},
  source={book/grade10/term_2/1-practice_activities}
}

\HighlightDeclareConcept{design-count}{
  role={count.design},
  notation={$\HLObservationCount{n}$},
  name={\HLObservationCount{Observation} \HLObservationCount{count}},
  brief={Fixed trial count or sample size; a positive integer.},
  definition={$n$ is the fixed number of trials or selected units; it is a positive integer. It differs from the size $N$ of the whole population. A subscript, as in $n_i$, identifies the sample.},
  source={book/grade10/term_2/1-practice_activities}
}

\HighlightDeclareConcept{wait}{
  role={quantity.primary},
  notation={$\GTenHLValue{T}$, $\GTenHLValue{t}$; $\GTenHLValue{F}=\GTenHLValue{T}-1$ or $\GTenHLValue{F}=\GTenHLValue{T}-\HLModelControl{r}$},
  name={\GTenHLValue{Trials} \GTenHLValue{and} \GTenHLValue{failures}},
  brief={Total trials include the final success; failures exclude successes.},
  definition={$T$ counts trials until the target is reached, including the final success; $t$ is a possible value. $F=T-1$ counts failures before the first success and $F=T-r$ before the $r$th success.},
  source={book/grade10/term_2/1-practice_activities}
}

\HighlightDeclareConcept{target}{
  role={model.control},
  notation={$\HLModelControl{r}$},
  name={\HLModelControl{Success} \HLModelControl{target}},
  brief={Fixed number of successes required before stopping.},
  definition={$r$ is the positive integer number of successes required to stop the experiment. It is a fixed model condition, rather than the random number of trials needed.},
  source={book/grade10/term_2/1-practice_activities}
}

\HighlightDeclareConcept{event-rate-per-exposure}{
  role={context},
  notation={$\mathrm{rate}$, $\mathrm{exposure}$},
  name={\HLMathContext{Rate} \HLMathContext{and} \HLMathContext{exposure}},
  brief={Events per exposure unit and length of the observation window.},
  definition={$\mathrm{rate}$ is the mean number of events per unit of exposure; $\mathrm{exposure}$ is the observed duration or extent. Multiply them using compatible units to obtain $\lambda$.},
  source={book/grade10/term_2/1-practice_activities}
}

\HighlightDeclareConcept{secondmoment}{
  role={expectation},
  notation={$\HLExpectedValue{E}(\GTenHLValue{X^2})$, $[\HLExpectedValue{E}(\GTenHLValue{X})]^2$},
  name={\HLExpectedValue{Second} \HLExpectedValue{moment} \HLExpectedValue{and} \HLExpectedValue{squared} \HLExpectedValue{mean}},
  brief={Average of squares; square of the average. Both use squared units.},
  definition={$E(X^2)$ averages the squares; $[E(X)]^2$ squares the average. These are distinct quantities with squared units; their difference is $\operatorname{Var}(X)$, if the second moment exists.},
  source={book/grade10/term_2/1-practice_activities}
}

\HighlightDeclareContext{g10-t2-c1-c2-c3-c4-c5-c6-c7-c8-c9-10-outcome}{
  concept={outcome-space},
  grade={10},
  term={2},
  criteria={1,2,3,4,5,6,7,8,9},
  variant={legacy},
  notation={$\GTenHLOutcome{\omega}$, $\GTenHLOutcome{\Omega}$},
  examples={},
  source={book/grade10/term_2/1-practice_activities/G10_T2_L1_C1_practice_activity.tex; book/grade10/term_2/1-practice_activities/G10_T2_L1_C2_practice_activity.tex; book/grade10/term_2/1-practice_activities/G10_T2_L1_C3_practice_activity.tex; book/grade10/term_2/1-practice_activities/G10_T2_L1_C4_practice_activity.tex; book/grade10/term_2/1-practice_activities/G10_T2_L2_C5_practice_activity.tex; book/grade10/term_2/1-practice_activities/G10_T2_L2_C6_practice_activity.tex; book/grade10/term_2/1-practice_activities/G10_T2_L2_C7_practice_activity.tex; book/grade10/term_2/1-practice_activities/G10_T2_L3_C8_practice_activity.tex; book/grade10/term_2/1-practice_activities/G10_T2_L3_C9_practice_activity.tex},
  brief={One complete result; all possible complete results.}
}

\HighlightDeclareContext{g10-t2-c1-c2-c3-c4-c5-c6-c7-c8-c9-10-mass}{
  concept={event-probability},
  grade={10},
  term={2},
  criteria={1,2,3,4,5,6,7,8,9},
  variant={legacy},
  notation={$\HLProbability{P(A)}$, $\HLProbability{P(\omega)}$, $\HLProbability{p_X(x)}$, $\HLProbability{m_X(x)}$},
  examples={},
  source={book/grade10/term_2/1-practice_activities/G10_T2_L1_C1_practice_activity.tex; book/grade10/term_2/1-practice_activities/G10_T2_L1_C2_practice_activity.tex; book/grade10/term_2/1-practice_activities/G10_T2_L1_C3_practice_activity.tex; book/grade10/term_2/1-practice_activities/G10_T2_L1_C4_practice_activity.tex; book/grade10/term_2/1-practice_activities/G10_T2_L2_C5_practice_activity.tex; book/grade10/term_2/1-practice_activities/G10_T2_L2_C6_practice_activity.tex; book/grade10/term_2/1-practice_activities/G10_T2_L2_C7_practice_activity.tex; book/grade10/term_2/1-practice_activities/G10_T2_L3_C8_practice_activity.tex; book/grade10/term_2/1-practice_activities/G10_T2_L3_C9_practice_activity.tex},
  brief={Unitless chance of an event or one supported value.},
  definition={$P(A)$ is the probability of an event; $P(\omega)$ abbreviates $P(\{\omega\})$. $p_X(x)=m_X(x)=P(X=x)$ is the mass at $x$; $P(X=k)$ and $P(T=t)$ also occur. These are unitless quantities between 0 and 1.}
}

\HighlightDeclareContext{g10-t2-c1-c2-c3-c4-c5-c6-c7-c8-c9-10-variable}{
  concept={random-variable},
  grade={10},
  term={2},
  criteria={1,2,3,4,5,6,7,8,9},
  variant={legacy},
  notation={$\GTenHLValue{X}$, $\GTenHLValue{x}$, $\GTenHLValue{T}$, $\GTenHLValue{t}$; $\GTenHLValue{X}(\GTenHLOutcome{\omega})$},
  examples={},
  source={book/grade10/term_2/1-practice_activities/G10_T2_L1_C1_practice_activity.tex; book/grade10/term_2/1-practice_activities/G10_T2_L1_C2_practice_activity.tex; book/grade10/term_2/1-practice_activities/G10_T2_L1_C3_practice_activity.tex; book/grade10/term_2/1-practice_activities/G10_T2_L1_C4_practice_activity.tex; book/grade10/term_2/1-practice_activities/G10_T2_L2_C5_practice_activity.tex; book/grade10/term_2/1-practice_activities/G10_T2_L2_C6_practice_activity.tex; book/grade10/term_2/1-practice_activities/G10_T2_L2_C7_practice_activity.tex; book/grade10/term_2/1-practice_activities/G10_T2_L3_C8_practice_activity.tex; book/grade10/term_2/1-practice_activities/G10_T2_L3_C9_practice_activity.tex},
  brief={Numerical rule and possible value; state the units.},
  definition={$X$ or $T$ assigns a number to each outcome; $x$ or $t$ is a possible value. $X(\omega)$ is the value at outcome $\omega$. Units are retained; $F$ can also name a primary quantity.}
}

\HighlightDeclareContext{g10-t2-c1-c2-c3-c4-c5-c6-c7-c8-c9-10-second}{
  concept={second},
  grade={10},
  term={2},
  criteria={1,2,3,4,5,6,7,8,9},
  variant={legacy},
  notation={$\GTenHLSecond{Y}$, $\GTenHLSecond{y}$; $\GTenHLSecond{Y}=f(\GTenHLValue{X})$; $\GTenHLSecond{Y_i}=f_i(\GTenHLValue{X})$},
  examples={},
  source={book/grade10/term_2/1-practice_activities/G10_T2_L1_C1_practice_activity.tex; book/grade10/term_2/1-practice_activities/G10_T2_L1_C2_practice_activity.tex; book/grade10/term_2/1-practice_activities/G10_T2_L1_C3_practice_activity.tex; book/grade10/term_2/1-practice_activities/G10_T2_L1_C4_practice_activity.tex; book/grade10/term_2/1-practice_activities/G10_T2_L2_C5_practice_activity.tex; book/grade10/term_2/1-practice_activities/G10_T2_L2_C6_practice_activity.tex; book/grade10/term_2/1-practice_activities/G10_T2_L2_C7_practice_activity.tex; book/grade10/term_2/1-practice_activities/G10_T2_L3_C8_practice_activity.tex; book/grade10/term_2/1-practice_activities/G10_T2_L3_C9_practice_activity.tex},
  brief={Another numerical description, with its own units.}
}

\HighlightDeclareContext{g10-t2-c1-c2-c3-c4-c5-c6-c7-c8-c9-10-third}{
  concept={third},
  grade={10},
  term={2},
  criteria={1,2,3,4,5,6,7,8,9},
  variant={legacy},
  notation={$\GTenHLThird{Z}$, $\GTenHLThird{z}$, $\GTenHLThird{C}$},
  examples={},
  source={book/grade10/term_2/1-practice_activities/G10_T2_L1_C1_practice_activity.tex; book/grade10/term_2/1-practice_activities/G10_T2_L1_C2_practice_activity.tex; book/grade10/term_2/1-practice_activities/G10_T2_L1_C3_practice_activity.tex; book/grade10/term_2/1-practice_activities/G10_T2_L1_C4_practice_activity.tex; book/grade10/term_2/1-practice_activities/G10_T2_L2_C5_practice_activity.tex; book/grade10/term_2/1-practice_activities/G10_T2_L2_C6_practice_activity.tex; book/grade10/term_2/1-practice_activities/G10_T2_L2_C7_practice_activity.tex; book/grade10/term_2/1-practice_activities/G10_T2_L3_C8_practice_activity.tex; book/grade10/term_2/1-practice_activities/G10_T2_L3_C9_practice_activity.tex},
  brief={A third numerical description; an indicator is unitless.}
}

\HighlightDeclareContext{g10-t2-c1-c2-c3-c4-c5-c6-c7-c8-c9-10-unknown}{
  concept={unknown-constant},
  grade={10},
  term={2},
  criteria={1,2,3,4,5,6,7,8,9},
  variant={legacy},
  notation={$\GTenHLParameter{k}$, $\GTenHLParameter{q}$, $\GTenHLParameter{c}$, $\GTenHLParameter{a}$, $\GTenHLParameter{w}$},
  examples={},
  source={book/grade10/term_2/1-practice_activities/G10_T2_L1_C1_practice_activity.tex; book/grade10/term_2/1-practice_activities/G10_T2_L1_C2_practice_activity.tex; book/grade10/term_2/1-practice_activities/G10_T2_L1_C3_practice_activity.tex; book/grade10/term_2/1-practice_activities/G10_T2_L1_C4_practice_activity.tex; book/grade10/term_2/1-practice_activities/G10_T2_L2_C5_practice_activity.tex; book/grade10/term_2/1-practice_activities/G10_T2_L2_C6_practice_activity.tex; book/grade10/term_2/1-practice_activities/G10_T2_L2_C7_practice_activity.tex; book/grade10/term_2/1-practice_activities/G10_T2_L3_C8_practice_activity.tex; book/grade10/term_2/1-practice_activities/G10_T2_L3_C9_practice_activity.tex},
  brief={A model quantity to determine, with a context-defined domain.},
  definition={$k,q,c,a,w$ are quantities to be determined from the model's constraints. Context sets their domains and units; $q$ here does not automatically mean $1-p$.}
}

\HighlightDeclareContext{g10-t2-c8-10-variable}{
  concept={random-variable},
  grade={10},
  term={2},
  criteria={8},
  variant={legacy},
  notation={$\GTenHLValue{X}$, $\GTenHLValue{x}$, $\GTenHLValue{T}$, $\GTenHLValue{t}$; $\GTenHLValue{X}(\GTenHLOutcome{\omega})$},
  examples={},
  source={book/grade10/term_2/1-practice_activities/G10_T2_L3_C8_practice_activity.tex},
  brief={Numerical rule and possible value; state the units.},
  definition={$X$ or $T$ assigns a number to each outcome; $x$ or $t$ is a possible value. $X(\omega)$ is the value at outcome $\omega$. Units are retained; $F$ can also name a primary quantity.}
}

\HighlightDeclareContext{g10-t2-c8-10-second}{
  concept={second},
  grade={10},
  term={2},
  criteria={8},
  variant={legacy},
  notation={$\GTenHLSecond{Y}$, $\GTenHLSecond{y}$; $\GTenHLSecond{Y}=f(\GTenHLValue{X})$; $\GTenHLSecond{Y_i}=f_i(\GTenHLValue{X})$},
  examples={},
  source={book/grade10/term_2/1-practice_activities/G10_T2_L3_C8_practice_activity.tex},
  brief={Another numerical description, with its own units.}
}

\HighlightDeclareContext{g10-t2-c8-10-expected}{
  concept={expected-value},
  grade={10},
  term={2},
  criteria={8},
  variant={legacy},
  notation={$\HLExpectedValue{E}(\GTenHLValue{X})$, $\HLExpectedValue{\mathbb E}[\GTenHLSecond{Y_i}]$},
  examples={},
  source={book/grade10/term_2/1-practice_activities/G10_T2_L3_C8_practice_activity.tex},
  brief={Probability-weighted average, in original units.},
  definition={$E(X)$ and $\mathbb E[Y_i]$ are probability-weighted averages, when they exist. They have the variable's units; $E(Y)$ is also used. The operator $E$ differs from an event or indicator named $E$.}
}

\HighlightDeclareContext{g10-t2-c8-10-secondmoment}{
  concept={secondmoment},
  grade={10},
  term={2},
  criteria={8},
  variant={legacy},
  notation={$\HLExpectedValue{E}(\GTenHLValue{X^2})$, $[\HLExpectedValue{E}(\GTenHLValue{X})]^2$},
  examples={},
  source={book/grade10/term_2/1-practice_activities/G10_T2_L3_C8_practice_activity.tex},
  brief={Average of squares; square of the average. Both use squared units.}
}

\HighlightDeclareContext{g10-t2-c1-coordinated-10-outcome}{
  concept={outcome-space},
  grade={10},
  term={2},
  criteria={1},
  variant={coordinated},
  notation={$\GTenHLOutcome{\omega}$, $\GTenHLOutcome{\Omega}$},
  examples={},
  source={book/grade10/term_2/1-practice_activities/G10_T2_L1_C1_practice_activity.tex},
  brief={One complete result; all possible complete results.}
}

\HighlightDeclareContext{g10-t2-c1-coordinated-10-event}{
  concept={event},
  grade={10},
  term={2},
  criteria={1},
  variant={coordinated},
  notation={$\HLMathEvent{A}$},
  examples={},
  source={book/grade10/term_2/1-practice_activities/G10_T2_L1_C1_practice_activity.tex},
  brief={Outcomes meeting a condition; those outside it.}
}

\HighlightDeclareContext{g10-t2-c1-coordinated-10-variable}{
  concept={random-variable},
  grade={10},
  term={2},
  criteria={1},
  variant={coordinated},
  notation={$\GTenHLValue{X}$, $\GTenHLValue{x}$, $\GTenHLValue{F}$; $\GTenHLValue{X}(\GTenHLOutcome{\omega})$},
  examples={},
  source={book/grade10/term_2/1-practice_activities/G10_T2_L1_C1_practice_activity.tex},
  brief={Numerical rule and possible value; state the units.},
  definition={$X$ or $T$ assigns a number to each outcome; $x$ or $t$ is a possible value. $X(\omega)$ is the value at outcome $\omega$. Units are retained; $F$ can also name a primary quantity.}
}

\HighlightDeclareContext{g10-t2-c1-coordinated-10-support}{
  concept={discrete-support},
  grade={10},
  term={2},
  criteria={1},
  variant={coordinated},
  notation={$\GTenHLValue{S_X}=\{\GTenHLValue{x}:\HLProbability{P}(\GTenHLValue{X}=\GTenHLValue{x})>0\}$},
  examples={},
  source={book/grade10/term_2/1-practice_activities/G10_T2_L1_C1_practice_activity.tex},
  brief={Distinct values with positive probability, in the variable's units.}
}

\HighlightDeclareContext{g10-t2-c1-coordinated-10-second}{
  concept={second},
  grade={10},
  term={2},
  criteria={1},
  variant={coordinated},
  notation={$\GTenHLSecond{Y}$, $\GTenHLSecond{y}$, $\GTenHLSecond{E}$},
  examples={},
  source={book/grade10/term_2/1-practice_activities/G10_T2_L1_C1_practice_activity.tex},
  brief={Another numerical description, with its own units.}
}

\HighlightDeclareContext{g10-t2-c1-coordinated-10-third}{
  concept={third},
  grade={10},
  term={2},
  criteria={1},
  variant={coordinated},
  notation={$\GTenHLThird{Z}$, $\GTenHLThird{z}$, $\GTenHLThird{C}$},
  examples={},
  source={book/grade10/term_2/1-practice_activities/G10_T2_L1_C1_practice_activity.tex},
  brief={A third numerical description; an indicator is unitless.}
}

\HighlightDeclareContext{g10-t2-c1-coordinated-10-mass}{
  concept={event-probability},
  grade={10},
  term={2},
  criteria={1},
  variant={coordinated},
  notation={$\HLProbability{P(\omega)}$, $\HLProbability{p_X(x)}$},
  examples={},
  source={book/grade10/term_2/1-practice_activities/G10_T2_L1_C1_practice_activity.tex},
  brief={Unitless chance of an event or one supported value.},
  definition={$P(A)$ is the probability of an event; $P(\omega)$ abbreviates $P(\{\omega\})$. $p_X(x)=m_X(x)=P(X=x)$ is the mass at $x$; $P(X=k)$ and $P(T=t)$ also occur. These are unitless quantities between 0 and 1.}
}

\HighlightDeclareContext{g10-t2-c1-coordinated-10-unknown}{
  concept={unknown-constant},
  grade={10},
  term={2},
  criteria={1},
  variant={coordinated},
  notation={$\GTenHLParameter{k}$, $\GTenHLParameter{q}$, $\GTenHLParameter{c}$},
  examples={},
  source={book/grade10/term_2/1-practice_activities/G10_T2_L1_C1_practice_activity.tex},
  brief={A model quantity to determine, with a context-defined domain.},
  definition={$k,q,c,a,w$ are quantities to be determined from the model's constraints. Context sets their domains and units; $q$ here does not automatically mean $1-p$.}
}

\HighlightDeclareContext{g10-t2-c1-coordinated-10-ops}{
  concept={mathematical-operations},
  grade={10},
  term={2},
  criteria={1},
  variant={coordinated},
  notation={$\sum$},
  examples={},
  source={book/grade10/term_2/1-practice_activities/G10_T2_L1_C1_practice_activity.tex},
  brief={Sum, count selections, take factorials, use the exponential constant, or name a model.}
}

\HighlightDeclareContext{g10-t2-c1-coordinated-10-assume}{
  concept={assume},
  grade={10},
  term={2},
  criteria={1},
  variant={coordinated},
  notation={Independence, constant $p$, fixed observation rule},
  examples={},
  source={book/grade10/term_2/1-practice_activities/G10_T2_L1_C1_practice_activity.tex},
  brief={Conditions to check before using the selected distribution.}
}

\HighlightDeclareContext{g10-t2-c2-coordinated-10-outcome}{
  concept={outcome-space},
  grade={10},
  term={2},
  criteria={2},
  variant={coordinated},
  notation={$\GTenHLOutcome{\omega}$, $\GTenHLOutcome{\Omega}$},
  examples={},
  source={book/grade10/term_2/1-practice_activities/G10_T2_L1_C2_practice_activity.tex},
  brief={One complete result; all possible complete results.}
}

\HighlightDeclareContext{g10-t2-c2-coordinated-10-event}{
  concept={event},
  grade={10},
  term={2},
  criteria={2},
  variant={coordinated},
  notation={$\{\GTenHLValue{X}\leq\GTenHLValue{a}\}$},
  examples={},
  source={book/grade10/term_2/1-practice_activities/G10_T2_L1_C2_practice_activity.tex},
  brief={Outcomes meeting a condition; those outside it.}
}

\HighlightDeclareContext{g10-t2-c2-coordinated-10-variable}{
  concept={random-variable},
  grade={10},
  term={2},
  criteria={2},
  variant={coordinated},
  notation={$\GTenHLValue{X}$, $\GTenHLValue{x}$},
  examples={},
  source={book/grade10/term_2/1-practice_activities/G10_T2_L1_C2_practice_activity.tex},
  brief={Numerical rule and possible value; state the units.},
  definition={$X$ or $T$ assigns a number to each outcome; $x$ or $t$ is a possible value. $X(\omega)$ is the value at outcome $\omega$. Units are retained; $F$ can also name a primary quantity.}
}

\HighlightDeclareContext{g10-t2-c2-coordinated-10-support}{
  concept={discrete-support},
  grade={10},
  term={2},
  criteria={2},
  variant={coordinated},
  notation={$\GTenHLValue{S_X}=\{\GTenHLValue{x}:\HLProbability{P}(\GTenHLValue{X}=\GTenHLValue{x})>0\}$},
  examples={},
  source={book/grade10/term_2/1-practice_activities/G10_T2_L1_C2_practice_activity.tex},
  brief={Distinct values with positive probability, in the variable's units.}
}

\HighlightDeclareContext{g10-t2-c2-coordinated-10-mass}{
  concept={event-probability},
  grade={10},
  term={2},
  criteria={2},
  variant={coordinated},
  notation={$\HLProbability{p_X(x)}$},
  examples={},
  source={book/grade10/term_2/1-practice_activities/G10_T2_L1_C2_practice_activity.tex},
  brief={Unitless chance of an event or one supported value.},
  definition={$P(A)$ is the probability of an event; $P(\omega)$ abbreviates $P(\{\omega\})$. $p_X(x)=m_X(x)=P(X=x)$ is the mass at $x$; $P(X=k)$ and $P(T=t)$ also occur. These are unitless quantities between 0 and 1.}
}

\HighlightDeclareContext{g10-t2-c2-coordinated-10-cum}{
  concept={cumulative-probability},
  grade={10},
  term={2},
  criteria={2},
  variant={coordinated},
  notation={$\HLCumulativeProbability{F_X(a)}$, $\HLCumulativeProbability{F_X(a^-)}$; $\HLCumulativeProbability{P}(\GTenHLValue{X}\leq \GTenHLValue{a})$},
  examples={},
  source={book/grade10/term_2/1-practice_activities/G10_T2_L1_C2_practice_activity.tex},
  brief={Chance up to a cutoff; the left limit excludes the cutoff.},
  definition={$F_X(a)=P(X\leq a)$ includes the cutoff; $F_X(a^-)=P(X<a)$ excludes it. For a discrete variable, $F_X(a)-F_X(a^-)=p_X(a)=P(X=a)$: the jump is the point mass, including zero at an unsupported cutoff. $P(T\leq t)$ is also used. A cumulative height is not an individual probability mass.}
}

\HighlightDeclareContext{g10-t2-c2-coordinated-10-unknown}{
  concept={unknown-constant},
  grade={10},
  term={2},
  criteria={2},
  variant={coordinated},
  notation={$\GTenHLParameter{k}$, $\GTenHLParameter{q}$, $\GTenHLParameter{c}$, $\GTenHLParameter{a}$, $\GTenHLParameter{w}$},
  examples={},
  source={book/grade10/term_2/1-practice_activities/G10_T2_L1_C2_practice_activity.tex},
  brief={A model quantity to determine, with a context-defined domain.},
  definition={$k,q,c,a,w$ are quantities to be determined from the model's constraints. Context sets their domains and units; $q$ here does not automatically mean $1-p$.}
}

\HighlightDeclareContext{g10-t2-c2-coordinated-10-cutoff}{
  concept={cutoff},
  grade={10},
  term={2},
  criteria={2},
  variant={coordinated},
  notation={$\GTenHLValue{a}$, $\GTenHLValue{b}$, $\GTenHLValue{m}$, $\GTenHLValue{\lfloor x\rfloor}$},
  examples={},
  source={book/grade10/term_2/1-practice_activities/G10_T2_L1_C2_practice_activity.tex},
  brief={A boundary in the variable's units; specify included endpoints.}
}

\HighlightDeclareContext{g10-t2-c2-coordinated-10-ops}{
  concept={mathematical-operations},
  grade={10},
  term={2},
  criteria={2},
  variant={coordinated},
  notation={$\sum$, $\binom nk$, $n!$},
  examples={},
  source={book/grade10/term_2/1-practice_activities/G10_T2_L1_C2_practice_activity.tex},
  brief={Sum, count selections, take factorials, use the exponential constant, or name a model.}
}

\HighlightDeclareContext{g10-t2-c2-coordinated-10-assume}{
  concept={assume},
  grade={10},
  term={2},
  criteria={2},
  variant={coordinated},
  notation={Independence, constant $p$, fixed observation rule},
  examples={},
  source={book/grade10/term_2/1-practice_activities/G10_T2_L1_C2_practice_activity.tex},
  brief={Conditions to check before using the selected distribution.}
}

\HighlightDeclareContext{g10-t2-c2-coordinated-10-aux}{
  concept={aux},
  grade={10},
  term={2},
  criteria={2},
  variant={coordinated},
  notation={$K,a,b,c$; $a_m,b_m,c_m$; $N_k,b_j,R_m$},
  examples={},
  source={book/grade10/term_2/1-practice_activities/G10_T2_L1_C2_practice_activity.tex},
  brief={Problem-defined counts, recurrence states, balances or survival quantities.}
}

\HighlightDeclareContext{g10-t2-c3-coordinated-10-event}{
  concept={event},
  grade={10},
  term={2},
  criteria={3},
  variant={coordinated},
  notation={$\HLMathEvent{E}$, $\HLMathEvent{E^c}$},
  examples={},
  source={book/grade10/term_2/1-practice_activities/G10_T2_L1_C3_practice_activity.tex},
  brief={Outcomes meeting a condition; those outside it.}
}

\HighlightDeclareContext{g10-t2-c3-coordinated-10-variable}{
  concept={random-variable},
  grade={10},
  term={2},
  criteria={3},
  variant={coordinated},
  notation={$\GTenHLValue{X}$, $\GTenHLValue{x}$},
  examples={},
  source={book/grade10/term_2/1-practice_activities/G10_T2_L1_C3_practice_activity.tex},
  brief={Numerical rule and possible value; state the units.},
  definition={$X$ or $T$ assigns a number to each outcome; $x$ or $t$ is a possible value. $X(\omega)$ is the value at outcome $\omega$. Units are retained; $F$ can also name a primary quantity.}
}

\HighlightDeclareContext{g10-t2-c3-coordinated-10-support}{
  concept={discrete-support},
  grade={10},
  term={2},
  criteria={3},
  variant={coordinated},
  notation={$\GTenHLValue{S_X}=\{\GTenHLValue{x}:\HLProbability{P}(\GTenHLValue{X}=\GTenHLValue{x})>0\}$},
  examples={},
  source={book/grade10/term_2/1-practice_activities/G10_T2_L1_C3_practice_activity.tex},
  brief={Distinct values with positive probability, in the variable's units.}
}

\HighlightDeclareContext{g10-t2-c3-coordinated-10-p}{
  concept={success-probability},
  grade={10},
  term={2},
  criteria={3},
  variant={coordinated},
  notation={$\HLSuccessProbability{p}=\HLSuccessProbability{P}(\GTenHLValue{X}=\GTenHLValue{1})$},
  examples={},
  source={book/grade10/term_2/1-practice_activities/G10_T2_L1_C3_practice_activity.tex},
  brief={Chance of the specified state labelled 1, between 0 and 1.},
  definition={$p=P(X=1)$ for an indicator: the probability of the state called success. It is unitless and lies between 0 and 1; waiting-time models require $p>0$.}
}

\HighlightDeclareContext{g10-t2-c3-coordinated-10-q}{
  concept={complement-probability},
  grade={10},
  term={2},
  criteria={3},
  variant={coordinated},
  notation={$\HLComplementProbability{1-p}$},
  examples={},
  source={book/grade10/term_2/1-practice_activities/G10_T2_L1_C3_practice_activity.tex},
  brief={Chance of the complementary state labelled 0.},
  definition={$1-p$ is the complementary fraction or probability. Grade 10 also uses $q=1-p$; Grade 11 uses $1-\hat p$ and $1-p_0$ for sample and planning shares. Distinguish $q$ from an unrelated unknown.}
}

\HighlightDeclareContext{g10-t2-c3-coordinated-10-second}{
  concept={second},
  grade={10},
  term={2},
  criteria={3},
  variant={coordinated},
  notation={$\GTenHLSecond{Y_i}=f_i(\GTenHLValue{X})$},
  examples={},
  source={book/grade10/term_2/1-practice_activities/G10_T2_L1_C3_practice_activity.tex},
  brief={Another numerical description, with its own units.}
}

\HighlightDeclareContext{g10-t2-c3-coordinated-10-mass}{
  concept={event-probability},
  grade={10},
  term={2},
  criteria={3},
  variant={coordinated},
  notation={$\HLProbability{m_X(x)}$, $\HLProbability{m_{Y_i}(y)}$},
  examples={},
  source={book/grade10/term_2/1-practice_activities/G10_T2_L1_C3_practice_activity.tex},
  brief={Unitless chance of an event or one supported value.},
  definition={$P(A)$ is the probability of an event; $P(\omega)$ abbreviates $P(\{\omega\})$. $p_X(x)=m_X(x)=P(X=x)$ is the mass at $x$; $P(X=k)$ and $P(T=t)$ also occur. These are unitless quantities between 0 and 1.}
}

\HighlightDeclareContext{g10-t2-c3-coordinated-10-cum}{
  concept={cumulative-probability},
  grade={10},
  term={2},
  criteria={3},
  variant={coordinated},
  notation={$\HLCumulativeProbability{F_X(t)}$, $\HLCumulativeProbability{F_{Y_i}(t)}$},
  examples={},
  source={book/grade10/term_2/1-practice_activities/G10_T2_L1_C3_practice_activity.tex},
  brief={Chance up to a cutoff; the left limit excludes the cutoff.},
  definition={$F_X(a)=P(X\leq a)$ includes the cutoff; $F_X(a^-)=P(X<a)$ excludes it. $P(T\leq t)$ is also used. A cumulative probability is not an individual probability mass.}
}

\HighlightDeclareContext{g10-t2-c3-coordinated-10-expected}{
  concept={expected-value},
  grade={10},
  term={2},
  criteria={3},
  variant={coordinated},
  notation={$\HLExpectedValue{\mathbb E}[\GTenHLSecond{Y_i}]$},
  examples={},
  source={book/grade10/term_2/1-practice_activities/G10_T2_L1_C3_practice_activity.tex},
  brief={Probability-weighted average, in original units.},
  definition={$E(X)$ and $\mathbb E[Y_i]$ are probability-weighted averages, when they exist. They have the variable's units; $E(Y)$ is also used. The operator $E$ differs from an event or indicator named $E$.}
}

\HighlightDeclareContext{g10-t2-c3-coordinated-10-break}{
  concept={break},
  grade={10},
  term={2},
  criteria={3},
  variant={coordinated},
  notation={$\HLSuccessProbability{p_i^*}$},
  examples={},
  source={book/grade10/term_2/1-practice_activities/G10_T2_L1_C3_practice_activity.tex},
  brief={Probability at which the specified expected net payoff is zero.}
}

\HighlightDeclareContext{g10-t2-c3-coordinated-10-function}{
  concept={transform},
  role={transformation.function.c3},
  grade={10},
  term={2},
  criteria={3},
  variant={coordinated},
  notation={$\HLGtenCThreeFunction{f},\HLGtenCThreeFunction{f_i}$},
  name={\HLGtenCThreeFunction{Transformation function}},
  brief={\HLGtenCThreeFunction{Maps indicator states to monetary outcomes.}},
  definition={The rule maps indicator states to monetary outcomes; evaluated state payoffs retain the monetary-consequence role.},
  examples={},
  source={book/grade10/term_2/1-practice_activities/G10_T2_L1_C3_practice_activity.tex}
}

\HighlightDeclareContext{g10-t2-c3-coordinated-10-transform}{
  concept={transform},
  grade={10},
  term={2},
  criteria={3},
  variant={coordinated},
  notation={$f_i(0),f_i(1)$},
  examples={},
  source={book/grade10/term_2/1-practice_activities/G10_T2_L1_C3_practice_activity.tex},
  brief={Scale, fixed addition or state payoffs, with required units.}
}

\HighlightDeclareContext{g10-t2-c4-coordinated-10-variable}{
  concept={random-variable},
  grade={10},
  term={2},
  criteria={4},
  variant={coordinated},
  notation={$\GTenHLValue{X}$, $\GTenHLValue{x}$},
  examples={},
  source={book/grade10/term_2/1-practice_activities/G10_T2_L1_C4_practice_activity.tex},
  brief={Numerical rule and possible value; state the units.},
  definition={$X$ or $T$ assigns a number to each outcome; $x$ or $t$ is a possible value. $X(\omega)$ is the value at outcome $\omega$. Units are retained; $F$ can also name a primary quantity.}
}

\HighlightDeclareContext{g10-t2-c4-coordinated-10-support}{
  concept={discrete-support},
  grade={10},
  term={2},
  criteria={4},
  variant={coordinated},
  notation={$\GTenHLValue{S_X}=\{\GTenHLValue{x}:\HLProbability{P}(\GTenHLValue{X}=\GTenHLValue{x})>0\}$},
  examples={},
  source={book/grade10/term_2/1-practice_activities/G10_T2_L1_C4_practice_activity.tex},
  brief={Distinct values with positive probability, in the variable's units.}
}

\HighlightDeclareContext{g10-t2-c4-coordinated-10-second}{
  concept={second},
  grade={10},
  term={2},
  criteria={4},
  variant={coordinated},
  notation={$\GTenHLSecond{Y}$},
  examples={},
  source={book/grade10/term_2/1-practice_activities/G10_T2_L1_C4_practice_activity.tex},
  brief={Another numerical description, with its own units.}
}

\HighlightDeclareContext{g10-t2-c4-coordinated-10-p}{
  concept={success-probability},
  grade={10},
  term={2},
  criteria={4},
  variant={coordinated},
  notation={$\HLSuccessProbability{p}$},
  examples={},
  source={book/grade10/term_2/1-practice_activities/G10_T2_L1_C4_practice_activity.tex},
  brief={Chance of the specified state labelled 1, between 0 and 1.},
  definition={$p=P(X=1)$ for an indicator: the probability of the state called success. It is unitless and lies between 0 and 1; waiting-time models require $p>0$.}
}

\HighlightDeclareContext{g10-t2-c4-coordinated-10-q}{
  concept={complement-probability},
  grade={10},
  term={2},
  criteria={4},
  variant={coordinated},
  notation={$\HLComplementProbability{1-p}$},
  examples={},
  source={book/grade10/term_2/1-practice_activities/G10_T2_L1_C4_practice_activity.tex},
  brief={Chance of the complementary state labelled 0.},
  definition={$1-p$ is the complementary fraction or probability. Grade 10 also uses $q=1-p$; Grade 11 uses $1-\hat p$ and $1-p_0$ for sample and planning shares. Distinguish $q$ from an unrelated unknown.}
}

\HighlightDeclareContext{g10-t2-c4-coordinated-10-size}{
  concept={design-count},
  grade={10},
  term={2},
  criteria={4},
  variant={coordinated},
  notation={$\HLObservationCount{n}$},
  examples={},
  source={book/grade10/term_2/1-practice_activities/G10_T2_L1_C4_practice_activity.tex},
  brief={Fixed trial count or sample size; a positive integer.}
}

\HighlightDeclareContext{g10-t2-c4-coordinated-10-mass}{
  concept={event-probability},
  grade={10},
  term={2},
  criteria={4},
  variant={coordinated},
  notation={$\HLProbability{P(X=k)}$},
  examples={},
  source={book/grade10/term_2/1-practice_activities/G10_T2_L1_C4_practice_activity.tex},
  brief={Unitless chance of an event or one supported value.},
  definition={$P(A)$ is the probability of an event; $P(\omega)$ abbreviates $P(\{\omega\})$. $p_X(x)=m_X(x)=P(X=x)$ is the mass at $x$; $P(X=k)$ and $P(T=t)$ also occur. These are unitless quantities between 0 and 1.}
}

\HighlightDeclareContext{g10-t2-c4-coordinated-10-cum}{
  concept={cumulative-probability},
  grade={10},
  term={2},
  criteria={4},
  variant={coordinated},
  notation={$\HLCumulativeProbability{P(X\leq b)}$},
  examples={},
  source={book/grade10/term_2/1-practice_activities/G10_T2_L1_C4_practice_activity.tex},
  brief={Chance up to a cutoff; the left limit excludes the cutoff.},
  definition={$F_X(a)=P(X\leq a)$ includes the cutoff; $F_X(a^-)=P(X<a)$ excludes it. $P(T\leq t)$ is also used. A cumulative probability is not an individual probability mass.}
}

\HighlightDeclareContext{g10-t2-c4-coordinated-10-cutoff}{
  concept={cutoff},
  grade={10},
  term={2},
  criteria={4},
  variant={coordinated},
  notation={$\GTenHLValue{a}$, $\GTenHLValue{b}$, $\GTenHLValue{k}$},
  examples={},
  source={book/grade10/term_2/1-practice_activities/G10_T2_L1_C4_practice_activity.tex},
  brief={A boundary in the variable's units; specify included endpoints.}
}

\HighlightDeclareContext{g10-t2-c4-coordinated-10-ops}{
  concept={mathematical-operations},
  grade={10},
  term={2},
  criteria={4},
  variant={coordinated},
  notation={$\sum$, $\binom nk$, $n!$, $\sim$},
  examples={},
  source={book/grade10/term_2/1-practice_activities/G10_T2_L1_C4_practice_activity.tex},
  brief={Sum, count selections, take factorials, use the exponential constant, or name a model.}
}

\HighlightDeclareContext{g10-t2-c4-coordinated-10-assume}{
  concept={assume},
  grade={10},
  term={2},
  criteria={4},
  variant={coordinated},
  notation={Independence, constant $p$, fixed observation rule},
  examples={},
  source={book/grade10/term_2/1-practice_activities/G10_T2_L1_C4_practice_activity.tex},
  brief={Conditions to check before using the selected distribution.}
}

\HighlightDeclareContext{g10-t2-c5-coordinated-10-wait}{
  concept={wait},
  grade={10},
  term={2},
  criteria={5},
  variant={coordinated},
  notation={$\GTenHLValue{T}$, $\GTenHLValue{t}$; $\GTenHLValue{F}=\GTenHLValue{T}-1$},
  examples={},
  source={book/grade10/term_2/1-practice_activities/G10_T2_L2_C5_practice_activity.tex},
  brief={Total trials include the final success; failures exclude successes.}
}

\HighlightDeclareContext{g10-t2-c5-coordinated-10-support}{
  concept={discrete-support},
  grade={10},
  term={2},
  criteria={5},
  variant={coordinated},
  notation={$\GTenHLValue{S_T}$},
  examples={},
  source={book/grade10/term_2/1-practice_activities/G10_T2_L2_C5_practice_activity.tex},
  brief={Distinct values with positive probability, in the variable's units.}
}

\HighlightDeclareContext{g10-t2-c5-coordinated-10-p}{
  concept={success-probability},
  grade={10},
  term={2},
  criteria={5},
  variant={coordinated},
  notation={$\HLSuccessProbability{p}$},
  examples={},
  source={book/grade10/term_2/1-practice_activities/G10_T2_L2_C5_practice_activity.tex},
  brief={Chance of the specified state labelled 1, between 0 and 1.},
  definition={$p=P(X=1)$ for an indicator: the probability of the state called success. It is unitless and lies between 0 and 1; waiting-time models require $p>0$.}
}

\HighlightDeclareContext{g10-t2-c5-coordinated-10-q}{
  concept={complement-probability},
  grade={10},
  term={2},
  criteria={5},
  variant={coordinated},
  notation={$\HLComplementProbability{1-p}$},
  examples={},
  source={book/grade10/term_2/1-practice_activities/G10_T2_L2_C5_practice_activity.tex},
  brief={Chance of the complementary state labelled 0.},
  definition={$1-p$ is the complementary fraction or probability. Grade 10 also uses $q=1-p$; Grade 11 uses $1-\hat p$ and $1-p_0$ for sample and planning shares. Distinguish $q$ from an unrelated unknown.}
}

\HighlightDeclareContext{g10-t2-c5-coordinated-10-second}{
  concept={second},
  grade={10},
  term={2},
  criteria={5},
  variant={coordinated},
  notation={$\GTenHLSecond{C}$},
  examples={},
  source={book/grade10/term_2/1-practice_activities/G10_T2_L2_C5_practice_activity.tex},
  brief={Another numerical description, with its own units.}
}

\HighlightDeclareContext{g10-t2-c5-coordinated-10-mass}{
  concept={event-probability},
  grade={10},
  term={2},
  criteria={5},
  variant={coordinated},
  notation={$\HLProbability{P(T=t)}$},
  examples={},
  source={book/grade10/term_2/1-practice_activities/G10_T2_L2_C5_practice_activity.tex},
  brief={Unitless chance of an event or one supported value.},
  definition={$P(A)$ is the probability of an event; $P(\omega)$ abbreviates $P(\{\omega\})$. $p_X(x)=m_X(x)=P(X=x)$ is the mass at $x$; $P(X=k)$ and $P(T=t)$ also occur. These are unitless quantities between 0 and 1.}
}

\HighlightDeclareContext{g10-t2-c5-coordinated-10-cum}{
  concept={cumulative-probability},
  grade={10},
  term={2},
  criteria={5},
  variant={coordinated},
  notation={$\HLCumulativeProbability{P(T\leq t)}$},
  examples={},
  source={book/grade10/term_2/1-practice_activities/G10_T2_L2_C5_practice_activity.tex},
  brief={Chance up to a cutoff; the left limit excludes the cutoff.},
  definition={$F_X(a)=P(X\leq a)$ includes the cutoff; $F_X(a^-)=P(X<a)$ excludes it. $P(T\leq t)$ is also used. A cumulative probability is not an individual probability mass.}
}

\HighlightDeclareContext{g10-t2-c5-coordinated-10-cutoff}{
  concept={cutoff},
  grade={10},
  term={2},
  criteria={5},
  variant={coordinated},
  notation={$\GTenHLValue{m}$},
  examples={},
  source={book/grade10/term_2/1-practice_activities/G10_T2_L2_C5_practice_activity.tex},
  brief={A boundary in the variable's units; specify included endpoints.}
}

\HighlightDeclareContext{g10-t2-c5-coordinated-10-event}{
  concept={event},
  grade={10},
  term={2},
  criteria={5},
  variant={coordinated},
  notation={$\HLMathEvent{A}$, $\HLMathEvent{B}$},
  examples={},
  source={book/grade10/term_2/1-practice_activities/G10_T2_L2_C5_practice_activity.tex},
  brief={Outcomes meeting a condition; those outside it.}
}

\HighlightDeclareContext{g10-t2-c5-coordinated-10-assume}{
  concept={assume},
  grade={10},
  term={2},
  criteria={5},
  variant={coordinated},
  notation={Independence, constant $p$, fixed observation rule},
  examples={},
  source={book/grade10/term_2/1-practice_activities/G10_T2_L2_C5_practice_activity.tex},
  brief={Conditions to check before using the selected distribution.}
}

\HighlightDeclareContext{g10-t2-c6-coordinated-10-wait}{
  concept={wait},
  grade={10},
  term={2},
  criteria={6},
  variant={coordinated},
  notation={$\GTenHLValue{T}$, $\GTenHLValue{t}$; $\GTenHLValue{F}=\GTenHLValue{T}-\HLModelControl{r}$},
  examples={},
  source={book/grade10/term_2/1-practice_activities/G10_T2_L2_C6_practice_activity.tex},
  brief={Total trials include the final success; failures exclude successes.}
}

\HighlightDeclareContext{g10-t2-c6-coordinated-10-support}{
  concept={discrete-support},
  grade={10},
  term={2},
  criteria={6},
  variant={coordinated},
  notation={$\GTenHLValue{S_T}$},
  examples={},
  source={book/grade10/term_2/1-practice_activities/G10_T2_L2_C6_practice_activity.tex},
  brief={Distinct values with positive probability, in the variable's units.}
}

\HighlightDeclareContext{g10-t2-c6-coordinated-10-p}{
  concept={success-probability},
  grade={10},
  term={2},
  criteria={6},
  variant={coordinated},
  notation={$\HLSuccessProbability{p}$},
  examples={},
  source={book/grade10/term_2/1-practice_activities/G10_T2_L2_C6_practice_activity.tex},
  brief={Chance of the specified state labelled 1, between 0 and 1.},
  definition={$p=P(X=1)$ for an indicator: the probability of the state called success. It is unitless and lies between 0 and 1; waiting-time models require $p>0$.}
}

\HighlightDeclareContext{g10-t2-c6-coordinated-10-q}{
  concept={complement-probability},
  grade={10},
  term={2},
  criteria={6},
  variant={coordinated},
  notation={$\HLComplementProbability{1-p}$},
  examples={},
  source={book/grade10/term_2/1-practice_activities/G10_T2_L2_C6_practice_activity.tex},
  brief={Chance of the complementary state labelled 0.},
  definition={$1-p$ is the complementary fraction or probability. Grade 10 also uses $q=1-p$; Grade 11 uses $1-\hat p$ and $1-p_0$ for sample and planning shares. Distinguish $q$ from an unrelated unknown.}
}

\HighlightDeclareContext{g10-t2-c6-coordinated-10-target}{
  concept={target},
  grade={10},
  term={2},
  criteria={6},
  variant={coordinated},
  notation={$\HLModelControl{r}$},
  examples={},
  source={book/grade10/term_2/1-practice_activities/G10_T2_L2_C6_practice_activity.tex},
  brief={Fixed number of successes required before stopping.}
}

\HighlightDeclareContext{g10-t2-c6-coordinated-10-second}{
  concept={second},
  grade={10},
  term={2},
  criteria={6},
  variant={coordinated},
  notation={$\GTenHLSecond{C}$},
  examples={},
  source={book/grade10/term_2/1-practice_activities/G10_T2_L2_C6_practice_activity.tex},
  brief={Another numerical description, with its own units.}
}

\HighlightDeclareContext{g10-t2-c6-coordinated-10-mass}{
  concept={event-probability},
  grade={10},
  term={2},
  criteria={6},
  variant={coordinated},
  notation={$\HLProbability{P(T=t)}$},
  examples={},
  source={book/grade10/term_2/1-practice_activities/G10_T2_L2_C6_practice_activity.tex},
  brief={Unitless chance of an event or one supported value.},
  definition={$P(A)$ is the probability of an event; $P(\omega)$ abbreviates $P(\{\omega\})$. $p_X(x)=m_X(x)=P(X=x)$ is the mass at $x$; $P(X=k)$ and $P(T=t)$ also occur. These are unitless quantities between 0 and 1.}
}

\HighlightDeclareContext{g10-t2-c6-coordinated-10-cum}{
  concept={cumulative-probability},
  grade={10},
  term={2},
  criteria={6},
  variant={coordinated},
  notation={$\HLCumulativeProbability{P(T\leq t)}$},
  examples={},
  source={book/grade10/term_2/1-practice_activities/G10_T2_L2_C6_practice_activity.tex},
  brief={Chance up to a cutoff; the left limit excludes the cutoff.},
  definition={$F_X(a)=P(X\leq a)$ includes the cutoff; $F_X(a^-)=P(X<a)$ excludes it. $P(T\leq t)$ is also used. A cumulative probability is not an individual probability mass.}
}

\HighlightDeclareContext{g10-t2-c6-coordinated-10-size}{
  concept={design-count},
  grade={10},
  term={2},
  criteria={6},
  variant={coordinated},
  notation={$\HLObservationCount{n}$},
  examples={},
  source={book/grade10/term_2/1-practice_activities/G10_T2_L2_C6_practice_activity.tex},
  brief={Fixed trial count or sample size; a positive integer.}
}

\HighlightDeclareContext{g10-t2-c6-coordinated-10-cutoff}{
  concept={cutoff},
  grade={10},
  term={2},
  criteria={6},
  variant={coordinated},
  notation={$\GTenHLValue{m}$},
  examples={},
  source={book/grade10/term_2/1-practice_activities/G10_T2_L2_C6_practice_activity.tex},
  brief={A boundary in the variable's units; specify included endpoints.}
}

\HighlightDeclareContext{g10-t2-c6-coordinated-10-ops}{
  concept={mathematical-operations},
  grade={10},
  term={2},
  criteria={6},
  variant={coordinated},
  notation={$\sum$, $\binom nk$, $n!$, $\sim$},
  examples={},
  source={book/grade10/term_2/1-practice_activities/G10_T2_L2_C6_practice_activity.tex},
  brief={Sum, count selections, take factorials, use the exponential constant, or name a model.}
}

\HighlightDeclareContext{g10-t2-c6-coordinated-10-assume}{
  concept={assume},
  grade={10},
  term={2},
  criteria={6},
  variant={coordinated},
  notation={Independence, constant $p$, fixed observation rule},
  examples={},
  source={book/grade10/term_2/1-practice_activities/G10_T2_L2_C6_practice_activity.tex},
  brief={Conditions to check before using the selected distribution.}
}

\HighlightDeclareContext{g10-t2-c7-coordinated-10-variable}{
  concept={random-variable},
  grade={10},
  term={2},
  criteria={7},
  variant={coordinated},
  notation={$\GTenHLValue{X}$, $\GTenHLValue{k}$},
  examples={},
  source={book/grade10/term_2/1-practice_activities/G10_T2_L2_C7_practice_activity.tex},
  brief={Numerical rule and possible value; state the units.},
  definition={$X$ or $T$ assigns a number to each outcome; $x$ or $t$ is a possible value. $X(\omega)$ is the value at outcome $\omega$. Units are retained; $F$ can also name a primary quantity.}
}

\HighlightDeclareContext{g10-t2-c7-coordinated-10-support}{
  concept={discrete-support},
  grade={10},
  term={2},
  criteria={7},
  variant={coordinated},
  notation={$\GTenHLValue{S_X}=\{\GTenHLValue{x}:\HLProbability{P}(\GTenHLValue{X}=\GTenHLValue{x})>0\}$},
  examples={},
  source={book/grade10/term_2/1-practice_activities/G10_T2_L2_C7_practice_activity.tex},
  brief={Distinct values with positive probability, in the variable's units.}
}

\HighlightDeclareContext{g10-t2-c7-coordinated-10-mass}{
  concept={event-probability},
  grade={10},
  term={2},
  criteria={7},
  variant={coordinated},
  notation={$\HLProbability{P(X=k)}$},
  examples={},
  source={book/grade10/term_2/1-practice_activities/G10_T2_L2_C7_practice_activity.tex},
  brief={Unitless chance of an event or one supported value.},
  definition={$P(A)$ is the probability of an event; $P(\omega)$ abbreviates $P(\{\omega\})$. $p_X(x)=m_X(x)=P(X=x)$ is the mass at $x$; $P(X=k)$ and $P(T=t)$ also occur. These are unitless quantities between 0 and 1.}
}

\HighlightDeclareContext{g10-t2-c7-coordinated-10-cum}{
  concept={cumulative-probability},
  grade={10},
  term={2},
  criteria={7},
  variant={coordinated},
  notation={$\HLCumulativeProbability{P(X\leq b)}$},
  examples={},
  source={book/grade10/term_2/1-practice_activities/G10_T2_L2_C7_practice_activity.tex},
  brief={Chance up to a cutoff; the left limit excludes the cutoff.},
  definition={$F_X(a)=P(X\leq a)$ includes the cutoff; $F_X(a^-)=P(X<a)$ excludes it. $P(T\leq t)$ is also used. A cumulative probability is not an individual probability mass.}
}

\HighlightDeclareContext{g10-t2-c7-coordinated-10-cutoff}{
  concept={cutoff},
  grade={10},
  term={2},
  criteria={7},
  variant={coordinated},
  notation={$\GTenHLValue{a}$, $\GTenHLValue{b}$, $\GTenHLValue{k}$},
  examples={},
  source={book/grade10/term_2/1-practice_activities/G10_T2_L2_C7_practice_activity.tex},
  brief={A boundary in the variable's units; specify included endpoints.}
}

\HighlightDeclareContext{g10-t2-c7-coordinated-10-rate}{
  concept={event-rate-per-exposure},
  grade={10},
  term={2},
  criteria={7},
  variant={coordinated},
  notation={$\mathrm{rate}$, $\mathrm{exposure}$},
  examples={},
  source={book/grade10/term_2/1-practice_activities/G10_T2_L2_C7_practice_activity.tex},
  brief={Events per exposure unit and length of the observation window.}
}

\HighlightDeclareContext{g10-t2-c7-coordinated-10-lambda}{
  concept={poisson-expected-count},
  grade={10},
  term={2},
  criteria={7},
  variant={coordinated},
  notation={$\HLPoissonMean{\lambda}$},
  examples={},
  source={book/grade10/term_2/1-practice_activities/G10_T2_L2_C7_practice_activity.tex},
  brief={Expected count for the selected exposure.},
  definition={$\lambda$ is the expected count for the chosen exposure: rate times exposure. In the Poisson model it is also the numerical value of the count's variance. It is not the rate of an exponential model.}
}

\HighlightDeclareContext{g10-t2-c7-coordinated-10-second}{
  concept={second},
  grade={10},
  term={2},
  criteria={7},
  variant={coordinated},
  notation={$\GTenHLSecond{L}$},
  examples={},
  source={book/grade10/term_2/1-practice_activities/G10_T2_L2_C7_practice_activity.tex},
  brief={Another numerical description, with its own units.}
}

\HighlightDeclareContext{g10-t2-c7-coordinated-10-ops}{
  concept={mathematical-operations},
  grade={10},
  term={2},
  criteria={7},
  variant={coordinated},
  notation={$\sum$, $k!$, $e$, $\sim$},
  examples={},
  source={book/grade10/term_2/1-practice_activities/G10_T2_L2_C7_practice_activity.tex},
  brief={Sum, count selections, take factorials, use the exponential constant, or name a model.}
}

\HighlightDeclareContext{g10-t2-c7-coordinated-10-assume}{
  concept={assume},
  grade={10},
  term={2},
  criteria={7},
  variant={coordinated},
  notation={Constant rate, independent interval counts, isolated events},
  examples={},
  source={book/grade10/term_2/1-practice_activities/G10_T2_L2_C7_practice_activity.tex},
  brief={Conditions to check before using the selected distribution.}
}

\HighlightDeclareContext{g10-t2-c8-coordinated-10-variable}{
  concept={random-variable},
  grade={10},
  term={2},
  criteria={8},
  variant={coordinated},
  notation={$\GTenHLValue{X}$, $\GTenHLValue{x}$},
  examples={},
  source={book/grade10/term_2/1-practice_activities/G10_T2_L3_C8_practice_activity.tex},
  brief={Numerical rule and possible value; state the units.},
  definition={$X$ or $T$ assigns a number to each outcome; $x$ or $t$ is a possible value. $X(\omega)$ is the value at outcome $\omega$. Units are retained; $F$ can also name a primary quantity.}
}

\HighlightDeclareContext{g10-t2-c8-coordinated-10-support}{
  concept={discrete-support},
  grade={10},
  term={2},
  criteria={8},
  variant={coordinated},
  notation={$\GTenHLValue{S_X}=\{\GTenHLValue{x}:\HLProbability{P}(\GTenHLValue{X}=\GTenHLValue{x})>0\}$},
  examples={},
  source={book/grade10/term_2/1-practice_activities/G10_T2_L3_C8_practice_activity.tex},
  brief={Distinct values with positive probability, in the variable's units.}
}

\HighlightDeclareContext{g10-t2-c8-coordinated-10-second}{
  concept={second},
  grade={10},
  term={2},
  criteria={8},
  variant={coordinated},
  notation={$\GTenHLSecond{Y}$},
  examples={},
  source={book/grade10/term_2/1-practice_activities/G10_T2_L3_C8_practice_activity.tex},
  brief={Another numerical description, with its own units.}
}

\HighlightDeclareContext{g10-t2-c8-coordinated-10-mass}{
  concept={event-probability},
  grade={10},
  term={2},
  criteria={8},
  variant={coordinated},
  notation={$\HLProbability{p_X(x)}$},
  examples={},
  source={book/grade10/term_2/1-practice_activities/G10_T2_L3_C8_practice_activity.tex},
  brief={Unitless chance of an event or one supported value.},
  definition={$P(A)$ is the probability of an event; $P(\omega)$ abbreviates $P(\{\omega\})$. $p_X(x)=m_X(x)=P(X=x)$ is the mass at $x$; $P(X=k)$ and $P(T=t)$ also occur. These are unitless quantities between 0 and 1.}
}

\HighlightDeclareContext{g10-t2-c8-coordinated-10-expected}{
  concept={expected-value},
  grade={10},
  term={2},
  criteria={8},
  variant={coordinated},
  notation={$\HLExpectedValue{E}(\GTenHLValue{X})$, $\HLExpectedValue{E}(\GTenHLSecond{Y})$},
  examples={},
  source={book/grade10/term_2/1-practice_activities/G10_T2_L3_C8_practice_activity.tex},
  brief={Probability-weighted average, in original units.},
  definition={$E(X)$ and $\mathbb E[Y_i]$ are probability-weighted averages, when they exist. They have the variable's units; $E(Y)$ is also used. The operator $E$ differs from an event or indicator named $E$.}
}

\HighlightDeclareContext{g10-t2-c8-coordinated-10-meanpop}{
  concept={population-mean},
  grade={10},
  term={2},
  criteria={8},
  variant={coordinated},
  notation={$\HLPopulationMean{\mu}$},
  examples={},
  source={book/grade10/term_2/1-practice_activities/G10_T2_L3_C8_practice_activity.tex},
  brief={Mean of the population or specified probability model, in original units.},
  definition={$\mu$ is the mean of the defined population or distribution; in a probability model it equals $E(X)$ when this exists. It has the variable's units and does not change when another sample is drawn from the same population.}
}

\HighlightDeclareContext{g10-t2-c8-coordinated-10-secondmoment}{
  concept={secondmoment},
  grade={10},
  term={2},
  criteria={8},
  variant={coordinated},
  notation={$\HLExpectedValue{E}(\GTenHLValue{X^2})$, $[\HLExpectedValue{E}(\GTenHLValue{X})]^2$},
  examples={},
  source={book/grade10/term_2/1-practice_activities/G10_T2_L3_C8_practice_activity.tex},
  brief={Average of squares; square of the average. Both use squared units.}
}

\HighlightDeclareContext{g10-t2-c8-coordinated-10-variance}{
  concept={population-variance},
  grade={10},
  term={2},
  criteria={8},
  variant={coordinated},
  notation={$\HLPopulationVariance{\operatorname{Var}}(\GTenHLValue{X})$, $\HLPopulationVariance{\sigma_X^2}$},
  examples={},
  source={book/grade10/term_2/1-practice_activities/G10_T2_L3_C8_practice_activity.tex},
  brief={Mean squared deviation from the mean; squared units.},
  definition={$\operatorname{Var}(X)=\sigma_X^2=E[(X-\mu)^2]$ measures squared dispersion about the mean. Also written $\sigma^2$, it has squared units and differs from the sample variance $s^2$.}
}

\HighlightDeclareContext{g10-t2-c8-coordinated-10-sd}{
  concept={population-standard-deviation},
  grade={10},
  term={2},
  criteria={8},
  variant={coordinated},
  notation={$\HLPopulationSD{\sigma_X}$, $\HLPopulationSD{\sigma}$},
  examples={},
  source={book/grade10/term_2/1-practice_activities/G10_T2_L3_C8_practice_activity.tex},
  brief={Spread of individual values, in original units.},
  definition={$\sigma_X$ or $\sigma$ is the nonnegative square root of the population or model variance. It describes the spread of individual values in the original units; it is not the standard error of a mean.}
}

\HighlightDeclareContext{g10-t2-c8-coordinated-10-transform}{
  concept={transform},
  grade={10},
  term={2},
  criteria={8},
  variant={coordinated},
  notation={$a,b$},
  examples={},
  source={book/grade10/term_2/1-practice_activities/G10_T2_L3_C8_practice_activity.tex},
  brief={Scale, fixed addition or state payoffs, with required units.}
}

\HighlightDeclareContext{g10-t2-c8-coordinated-10-ops}{
  concept={mathematical-operations},
  grade={10},
  term={2},
  criteria={8},
  variant={coordinated},
  notation={$\sum_x$},
  examples={},
  source={book/grade10/term_2/1-practice_activities/G10_T2_L3_C8_practice_activity.tex},
  brief={Sum, count selections, take factorials, use the exponential constant, or name a model.}
}

\HighlightDeclareContext{g10-t2-c9-coordinated-10-event}{
  concept={event},
  grade={10},
  term={2},
  criteria={9},
  variant={coordinated},
  notation={$\{\GTenHLValue{X}\geq\GTenHLValue{a}\}$},
  examples={},
  source={book/grade10/term_2/1-practice_activities/G10_T2_L3_C9_practice_activity.tex},
  brief={Outcomes meeting a condition; those outside it.}
}

\HighlightDeclareContext{g10-t2-c9-coordinated-10-variable}{
  concept={random-variable},
  grade={10},
  term={2},
  criteria={9},
  variant={coordinated},
  notation={$\GTenHLValue{X}$, $\GTenHLValue{x}$},
  examples={},
  source={book/grade10/term_2/1-practice_activities/G10_T2_L3_C9_practice_activity.tex},
  brief={Numerical rule and possible value; state the units.},
  definition={$X$ or $T$ assigns a number to each outcome; $x$ or $t$ is a possible value. $X(\omega)$ is the value at outcome $\omega$. Units are retained; $F$ can also name a primary quantity.}
}

\HighlightDeclareContext{g10-t2-c9-coordinated-10-support}{
  concept={discrete-support},
  grade={10},
  term={2},
  criteria={9},
  variant={coordinated},
  notation={$\GTenHLValue{S_X}=\{\GTenHLValue{x}:\HLProbability{P}(\GTenHLValue{X}=\GTenHLValue{x})>0\}$},
  examples={},
  source={book/grade10/term_2/1-practice_activities/G10_T2_L3_C9_practice_activity.tex},
  brief={Distinct values with positive probability, in the variable's units.}
}

\HighlightDeclareContext{g10-t2-c9-coordinated-10-wait}{
  concept={wait},
  grade={10},
  term={2},
  criteria={9},
  variant={coordinated},
  notation={$\GTenHLValue{T}$, $\GTenHLValue{t}$; $\GTenHLValue{F}=\GTenHLValue{T}-1$ or $\GTenHLValue{F}=\GTenHLValue{T}-\HLModelControl{r}$},
  examples={},
  source={book/grade10/term_2/1-practice_activities/G10_T2_L3_C9_practice_activity.tex},
  brief={Total trials include the final success; failures exclude successes.}
}

\HighlightDeclareContext{g10-t2-c9-coordinated-10-second}{
  concept={second},
  grade={10},
  term={2},
  criteria={9},
  variant={coordinated},
  notation={$\GTenHLSecond{L}$},
  examples={},
  source={book/grade10/term_2/1-practice_activities/G10_T2_L3_C9_practice_activity.tex},
  brief={Another numerical description, with its own units.}
}

\HighlightDeclareContext{g10-t2-c9-coordinated-10-p}{
  concept={success-probability},
  grade={10},
  term={2},
  criteria={9},
  variant={coordinated},
  notation={$\HLSuccessProbability{p}$},
  examples={},
  source={book/grade10/term_2/1-practice_activities/G10_T2_L3_C9_practice_activity.tex},
  brief={Chance of the specified state labelled 1, between 0 and 1.},
  definition={$p=P(X=1)$ for an indicator: the probability of the state called success. It is unitless and lies between 0 and 1; waiting-time models require $p>0$.}
}

\HighlightDeclareContext{g10-t2-c9-coordinated-10-q}{
  concept={complement-probability},
  grade={10},
  term={2},
  criteria={9},
  variant={coordinated},
  notation={$\HLComplementProbability{1-p}$},
  examples={},
  source={book/grade10/term_2/1-practice_activities/G10_T2_L3_C9_practice_activity.tex},
  brief={Chance of the complementary state labelled 0.},
  definition={$1-p$ is the complementary fraction or probability. Grade 10 also uses $q=1-p$; Grade 11 uses $1-\hat p$ and $1-p_0$ for sample and planning shares. Distinguish $q$ from an unrelated unknown.}
}

\HighlightDeclareContext{g10-t2-c9-coordinated-10-size}{
  concept={design-count},
  grade={10},
  term={2},
  criteria={9},
  variant={coordinated},
  notation={$\HLObservationCount{n}$},
  examples={},
  source={book/grade10/term_2/1-practice_activities/G10_T2_L3_C9_practice_activity.tex},
  brief={Fixed trial count or sample size; a positive integer.}
}

\HighlightDeclareContext{g10-t2-c9-coordinated-10-target}{
  concept={target},
  grade={10},
  term={2},
  criteria={9},
  variant={coordinated},
  notation={$\HLModelControl{r}$},
  examples={},
  source={book/grade10/term_2/1-practice_activities/G10_T2_L3_C9_practice_activity.tex},
  brief={Fixed number of successes required before stopping.}
}

\HighlightDeclareContext{g10-t2-c9-coordinated-10-mass}{
  concept={event-probability},
  grade={10},
  term={2},
  criteria={9},
  variant={coordinated},
  notation={$\HLProbability{P(X=k)}$, $\HLProbability{P(T=t)}$},
  examples={},
  source={book/grade10/term_2/1-practice_activities/G10_T2_L3_C9_practice_activity.tex},
  brief={Unitless chance of an event or one supported value.},
  definition={$P(A)$ is the probability of an event; $P(\omega)$ abbreviates $P(\{\omega\})$. $p_X(x)=m_X(x)=P(X=x)$ is the mass at $x$; $P(X=k)$ and $P(T=t)$ also occur. These are unitless quantities between 0 and 1.}
}

\HighlightDeclareContext{g10-t2-c9-coordinated-10-cum}{
  concept={cumulative-probability},
  grade={10},
  term={2},
  criteria={9},
  variant={coordinated},
  notation={$\HLCumulativeProbability{F_X(a)}$, $\HLCumulativeProbability{F_X(a^-)}$; $\HLCumulativeProbability{P}(\GTenHLValue{X}\leq \GTenHLValue{a})$},
  examples={},
  source={book/grade10/term_2/1-practice_activities/G10_T2_L3_C9_practice_activity.tex},
  brief={Chance up to a cutoff; the left limit excludes the cutoff.},
  definition={$F_X(a)=P(X\leq a)$ includes the cutoff; $F_X(a^-)=P(X<a)$ excludes it. $P(T\leq t)$ is also used. A cumulative probability is not an individual probability mass.}
}

\HighlightDeclareContext{g10-t2-c9-coordinated-10-rate}{
  concept={event-rate-per-exposure},
  grade={10},
  term={2},
  criteria={9},
  variant={coordinated},
  notation={$\mathrm{rate}$, $\mathrm{exposure}$},
  examples={},
  source={book/grade10/term_2/1-practice_activities/G10_T2_L3_C9_practice_activity.tex},
  brief={Events per exposure unit and length of the observation window.}
}

\HighlightDeclareContext{g10-t2-c9-coordinated-10-lambda}{
  concept={poisson-expected-count},
  grade={10},
  term={2},
  criteria={9},
  variant={coordinated},
  notation={$\HLPoissonMean{\lambda}$},
  examples={},
  source={book/grade10/term_2/1-practice_activities/G10_T2_L3_C9_practice_activity.tex},
  brief={Expected count for the selected exposure.},
  definition={$\lambda$ is the expected count for the chosen exposure: rate times exposure. In the Poisson model it is also the numerical value of the count's variance. It is not the rate of an exponential model.}
}

\HighlightDeclareContext{g10-t2-c9-coordinated-10-ops}{
  concept={mathematical-operations},
  grade={10},
  term={2},
  criteria={9},
  variant={coordinated},
  notation={$\sum$, $\binom nk$, $n!$, $e$, $\sim$},
  examples={},
  source={book/grade10/term_2/1-practice_activities/G10_T2_L3_C9_practice_activity.tex},
  brief={Sum, count selections, take factorials, use the exponential constant, or name a model.}
}

\HighlightDeclareContext{g10-t2-c9-coordinated-10-assume}{
  concept={assume},
  grade={10},
  term={2},
  criteria={9},
  variant={coordinated},
  notation={Independence, constant $p$, fixed observation rule},
  examples={},
  source={book/grade10/term_2/1-practice_activities/G10_T2_L3_C9_practice_activity.tex},
  brief={Conditions to check before using the selected distribution.}
}

\HighlightDeclareContext{g10-t2-c1-10-outcome}{
  concept={outcome-space},
  grade={10},
  term={2},
  criteria={1},
  variant={legacy},
  notation={$\GTenHLOutcome{\omega}$, $\GTenHLOutcome{\Omega}$},
  examples={},
  source={book/grade10/term_2/1-practice_activities/G10_T2_L1_C1_practice_activity.tex},
  brief={One complete result; all possible complete results.}
}

\HighlightDeclareContext{g10-t2-c1-10-mass}{
  concept={event-probability},
  grade={10},
  term={2},
  criteria={1},
  variant={legacy},
  notation={$\HLProbability{P(A)}$, $\HLProbability{P(\omega)}$, $\HLProbability{p_X(x)}$, $\HLProbability{m_X(x)}$},
  examples={},
  source={book/grade10/term_2/1-practice_activities/G10_T2_L1_C1_practice_activity.tex},
  brief={Unitless chance of an event or one supported value.}
}

\HighlightDeclareContext{g10-t2-c1-10-variable}{
  concept={random-variable},
  grade={10},
  term={2},
  criteria={1},
  variant={legacy},
  notation={$\GTenHLValue{X}$, $\GTenHLValue{x}$, $\GTenHLValue{T}$, $\GTenHLValue{t}$; $\GTenHLValue{X}(\GTenHLOutcome{\omega})$},
  examples={},
  source={book/grade10/term_2/1-practice_activities/G10_T2_L1_C1_practice_activity.tex},
  brief={Numerical rule and possible value; state the units.}
}

\HighlightDeclareContext{g10-t2-c1-10-second}{
  concept={second},
  grade={10},
  term={2},
  criteria={1},
  variant={legacy},
  notation={$\GTenHLSecond{Y}$, $\GTenHLSecond{y}$; $\GTenHLSecond{Y}=f(\GTenHLValue{X})$; $\GTenHLSecond{Y_i}=f_i(\GTenHLValue{X})$},
  examples={},
  source={book/grade10/term_2/1-practice_activities/G10_T2_L1_C1_practice_activity.tex},
  brief={Another numerical description, with its own units.}
}

\HighlightDeclareContext{g10-t2-c1-10-third}{
  concept={third},
  grade={10},
  term={2},
  criteria={1},
  variant={legacy},
  notation={$\GTenHLThird{Z}$, $\GTenHLThird{z}$, $\GTenHLThird{C}$},
  examples={},
  source={book/grade10/term_2/1-practice_activities/G10_T2_L1_C1_practice_activity.tex},
  brief={A third numerical description; an indicator is unitless.}
}

\HighlightDeclareContext{g10-t2-c1-10-unknown}{
  concept={unknown-constant},
  grade={10},
  term={2},
  criteria={1},
  variant={legacy},
  notation={$\GTenHLParameter{k}$, $\GTenHLParameter{q}$, $\GTenHLParameter{c}$, $\GTenHLParameter{a}$, $\GTenHLParameter{w}$},
  examples={},
  source={book/grade10/term_2/1-practice_activities/G10_T2_L1_C1_practice_activity.tex},
  brief={A model quantity to determine, with a context-defined domain.}
}

\HighlightDeclareContext{g10-t2-intro-outcome}{
  concept={outcome-space},
  grade={10},
  term={2},
  criteria={},
  variant={introductory-prerequisite},
  notation={$\GTenHLOutcome{\omega}$, $\GTenHLOutcome{\Omega}$},
  examples={},
  source={book/grade10/term_2/1-practice_activities},
  brief={One complete result; all possible complete results.}
}

\HighlightDeclareContext{g10-t2-intro-event}{
  concept={event},
  grade={10},
  term={2},
  criteria={},
  variant={introductory-prerequisite},
  notation={$\HLMathEvent{A}$, $\HLMathEvent{E}$, $\HLMathEvent{E^c}$; $\{\GTenHLValue{X}\leq \GTenHLValue{a}\}$; $\HLMathEvent{B}$; $\{\GTenHLValue{X}\geq\GTenHLValue{a}\}$},
  examples={},
  source={book/grade10/term_2/1-practice_activities},
  brief={Outcomes meeting a condition; those outside it.}
}

\HighlightDeclareContext{g10-t2-intro-variable}{
  concept={random-variable},
  grade={10},
  term={2},
  criteria={},
  variant={introductory-prerequisite},
  notation={$\GTenHLValue{X}$, $\GTenHLValue{x}$, $\GTenHLValue{T}$, $\GTenHLValue{t}$; $\GTenHLValue{X}(\GTenHLOutcome{\omega})$; $\GTenHLValue{F}$},
  examples={},
  source={book/grade10/term_2/1-practice_activities},
  brief={Numerical rule and possible value; state the units.},
  definition={$X$ or $T$ assigns a number to each outcome; $x$ or $t$ is a possible value. $X(\omega)$ is the value at outcome $\omega$. Units are retained; $F$ can also name a primary quantity.}
}

\HighlightDeclareContext{g10-t2-intro-support}{
  concept={discrete-support},
  grade={10},
  term={2},
  criteria={},
  variant={introductory-prerequisite},
  notation={$\GTenHLValue{S_X}=\{\GTenHLValue{x}:\HLProbability{P}(\GTenHLValue{X}=\GTenHLValue{x})>0\}$; $\GTenHLValue{S_T}$},
  examples={},
  source={book/grade10/term_2/1-practice_activities},
  brief={Distinct values with positive probability, in the variable's units.}
}

\HighlightDeclareContext{g10-t2-intro-second}{
  concept={second},
  grade={10},
  term={2},
  criteria={},
  variant={introductory-prerequisite},
  notation={$\GTenHLSecond{Y}$, $\GTenHLSecond{y}$; $\GTenHLSecond{Y}=f(\GTenHLValue{X})$; $\GTenHLSecond{Y_i}=f_i(\GTenHLValue{X})$; $\GTenHLSecond{E}$, $\GTenHLSecond{C}$, $\GTenHLSecond{L}$},
  examples={},
  source={book/grade10/term_2/1-practice_activities},
  brief={Another numerical description, with its own units.}
}

\HighlightDeclareContext{g10-t2-intro-third}{
  concept={third},
  grade={10},
  term={2},
  criteria={},
  variant={introductory-prerequisite},
  notation={$\GTenHLThird{Z}$, $\GTenHLThird{z}$, $\GTenHLThird{C}$},
  examples={},
  source={book/grade10/term_2/1-practice_activities},
  brief={A third numerical description; an indicator is unitless.}
}

\HighlightDeclareContext{g10-t2-intro-transform}{
  concept={transform},
  grade={10},
  term={2},
  criteria={},
  variant={introductory-prerequisite},
  notation={$a,b$; $f_i(0),f_i(1)$},
  examples={},
  source={book/grade10/term_2/1-practice_activities},
  brief={Scale, fixed addition or state payoffs, with required units.}
}

\HighlightDeclareContext{g10-t2-intro-cutoff}{
  concept={cutoff},
  grade={10},
  term={2},
  criteria={},
  variant={introductory-prerequisite},
  notation={$\GTenHLValue{a}$, $\GTenHLValue{b}$, $\GTenHLValue{m}$, $\GTenHLValue{\lfloor x\rfloor}$; $\GTenHLValue{k}$},
  examples={},
  source={book/grade10/term_2/1-practice_activities},
  brief={A boundary in the variable's units; specify included endpoints.}
}

\HighlightDeclareContext{g10-t2-intro-unknown}{
  concept={unknown-constant},
  grade={10},
  term={2},
  criteria={},
  variant={introductory-prerequisite},
  notation={$\GTenHLParameter{k}$, $\GTenHLParameter{q}$, $\GTenHLParameter{c}$, $\GTenHLParameter{a}$, $\GTenHLParameter{w}$},
  examples={},
  source={book/grade10/term_2/1-practice_activities},
  brief={A model quantity to determine, with a context-defined domain.},
  definition={$k,q,c,a,w$ are quantities to be determined from the model's constraints. Context sets their domains and units; $q$ here does not automatically mean $1-p$.}
}

\HighlightDeclareContext{g10-t2-intro-aux}{
  concept={aux},
  grade={10},
  term={2},
  criteria={},
  variant={introductory-prerequisite},
  notation={$K,a,b,c$; $a_m,b_m,c_m$; $N_k,b_j,R_m$},
  examples={},
  source={book/grade10/term_2/1-practice_activities},
  brief={Problem-defined counts, recurrence states, balances or survival quantities.}
}

\HighlightDeclareContext{g10-t2-intro-wait}{
  concept={wait},
  grade={10},
  term={2},
  criteria={},
  variant={introductory-prerequisite},
  notation={$\GTenHLValue{T}$, $\GTenHLValue{t}$; $\GTenHLValue{F}=\GTenHLValue{T}-1$ or $\GTenHLValue{F}=\GTenHLValue{T}-\HLModelControl{r}$},
  examples={},
  source={book/grade10/term_2/1-practice_activities},
  brief={Total trials include the final success; failures exclude successes.}
}

\HighlightDeclareContext{g10-t2-intro-ops}{
  concept={mathematical-operations},
  grade={10},
  term={2},
  criteria={},
  variant={introductory-prerequisite},
  notation={$\sum$, $\binom nk$, $n!$, $e$, $\sim$; $k!$, $\sum_x$},
  examples={},
  source={book/grade10/term_2/1-practice_activities},
  brief={Sum, count selections, take factorials, use the exponential constant, or name a model.}
}

\HighlightDeclareContext{g10-t2-intro-mass}{
  concept={event-probability},
  grade={10},
  term={2},
  criteria={},
  variant={introductory-prerequisite},
  notation={$\HLProbability{P(A)}$, $\HLProbability{P(\omega)}$, $\HLProbability{p_X(x)}$, $\HLProbability{m_X(x)}$; $\HLProbability{P(X=k)}$, $\HLProbability{P(T=t)}$},
  examples={},
  source={book/grade10/term_2/1-practice_activities},
  brief={Unitless chance of an event or one supported value.},
  definition={$P(A)$ is the probability of an event; $P(\omega)$ abbreviates $P(\{\omega\})$. $p_X(x)=m_X(x)=P(X=x)$ is the mass at $x$; $P(X=k)$ and $P(T=t)$ also occur. These are unitless quantities between 0 and 1.}
}

\HighlightDeclareContext{g10-t2-intro-cum}{
  concept={cumulative-probability},
  grade={10},
  term={2},
  criteria={},
  variant={introductory-prerequisite},
  notation={$\HLCumulativeProbability{F_X(a)}$, $\HLCumulativeProbability{F_X(a^-)}$; $\HLCumulativeProbability{P}(\GTenHLValue{X}\leq \GTenHLValue{a})$; $\HLCumulativeProbability{P(T\leq t)}$},
  examples={},
  source={book/grade10/term_2/1-practice_activities},
  brief={Chance up to a cutoff; the left limit excludes the cutoff.},
  definition={$F_X(a)=P(X\leq a)$ includes the cutoff; $F_X(a^-)=P(X<a)$ excludes it. $P(T\leq t)$ is also used. A cumulative probability is not an individual probability mass.}
}

\HighlightDeclareContext{g10-t2-intro-p}{
  concept={success-probability},
  grade={10},
  term={2},
  criteria={},
  variant={introductory-prerequisite},
  notation={$\HLSuccessProbability{p}=\HLSuccessProbability{P}(\GTenHLValue{X}=\GTenHLValue{1})$},
  examples={},
  source={book/grade10/term_2/1-practice_activities},
  brief={Chance of the specified state labelled 1, between 0 and 1.},
  definition={$p=P(X=1)$ for an indicator: the probability of the state called success. It is unitless and lies between 0 and 1; waiting-time models require $p>0$.}
}

\HighlightDeclareContext{g10-t2-intro-q}{
  concept={complement-probability},
  grade={10},
  term={2},
  criteria={},
  variant={introductory-prerequisite},
  notation={$\HLComplementProbability{1-p}$; $\HLComplementProbability{1-\hat p}$, $\HLComplementProbability{1-p_0}$},
  examples={},
  source={book/grade10/term_2/1-practice_activities},
  brief={Chance of the complementary state labelled 0.},
  definition={$1-p$ is the complementary fraction or probability. Grade 10 also uses $q=1-p$; Grade 11 uses $1-\hat p$ and $1-p_0$ for sample and planning shares. Distinguish $q$ from an unrelated unknown.}
}

\HighlightDeclareContext{g10-t2-intro-break}{
  concept={break},
  grade={10},
  term={2},
  criteria={},
  variant={introductory-prerequisite},
  notation={$\HLSuccessProbability{p_i^*}$},
  examples={},
  source={book/grade10/term_2/1-practice_activities},
  brief={Probability at which the specified expected net payoff is zero.}
}

\HighlightDeclareContext{g10-t2-intro-rate}{
  concept={event-rate-per-exposure},
  grade={10},
  term={2},
  criteria={},
  variant={introductory-prerequisite},
  notation={$\mathrm{rate}$, $\mathrm{exposure}$},
  examples={},
  source={book/grade10/term_2/1-practice_activities},
  brief={Events per exposure unit and length of the observation window.}
}

\HighlightDeclareContext{g10-t2-intro-lambda}{
  concept={poisson-expected-count},
  grade={10},
  term={2},
  criteria={},
  variant={introductory-prerequisite},
  notation={$\HLPoissonMean{\lambda}$},
  examples={},
  source={book/grade10/term_2/1-practice_activities},
  brief={Expected count for the selected exposure.},
  definition={$\lambda$ is the expected count for the chosen exposure: rate times exposure. In the Poisson model it is also the numerical value of the count's variance. It is not the rate of an exponential model.}
}

\HighlightDeclareContext{g10-t2-intro-expected}{
  concept={expected-value},
  grade={10},
  term={2},
  criteria={},
  variant={introductory-prerequisite},
  notation={$\HLExpectedValue{E}(\GTenHLValue{X})$, $\HLExpectedValue{\mathbb E}[\GTenHLSecond{Y_i}]$; $\HLExpectedValue{E}(\GTenHLSecond{Y})$},
  examples={},
  source={book/grade10/term_2/1-practice_activities},
  brief={Probability-weighted average, in original units.},
  definition={$E(X)$ and $\mathbb E[Y_i]$ are probability-weighted averages, when they exist. They have the variable's units; $E(Y)$ is also used. The operator $E$ differs from an event or indicator named $E$.}
}

\HighlightDeclareContext{g10-t2-intro-meanpop}{
  concept={population-mean},
  grade={10},
  term={2},
  criteria={},
  variant={introductory-prerequisite},
  notation={$\HLPopulationMean{\mu}$},
  examples={},
  source={book/grade10/term_2/1-practice_activities},
  brief={Mean of the population or specified probability model, in original units.},
  definition={$\mu$ is the mean of the defined population or distribution; in a probability model it equals $E(X)$ when this exists. It has the variable's units and does not change when another sample is drawn from the same population.}
}

\HighlightDeclareContext{g10-t2-intro-secondmoment}{
  concept={secondmoment},
  grade={10},
  term={2},
  criteria={},
  variant={introductory-prerequisite},
  notation={$\HLExpectedValue{E}(\GTenHLValue{X^2})$, $[\HLExpectedValue{E}(\GTenHLValue{X})]^2$},
  examples={},
  source={book/grade10/term_2/1-practice_activities},
  brief={Average of squares; square of the average. Both use squared units.}
}

\HighlightDeclareContext{g10-t2-intro-variance}{
  concept={population-variance},
  grade={10},
  term={2},
  criteria={},
  variant={introductory-prerequisite},
  notation={$\HLPopulationVariance{\operatorname{Var}}(\GTenHLValue{X})$, $\HLPopulationVariance{\sigma_X^2}$; $\HLPopulationVariance{\sigma^2}$},
  examples={},
  source={book/grade10/term_2/1-practice_activities},
  brief={Mean squared deviation from the mean; squared units.},
  definition={$\operatorname{Var}(X)=\sigma_X^2=E[(X-\mu)^2]$ measures squared dispersion about the mean. Also written $\sigma^2$, it has squared units and differs from the sample variance $s^2$.}
}

\HighlightDeclareContext{g10-t2-intro-sd}{
  concept={population-standard-deviation},
  grade={10},
  term={2},
  criteria={},
  variant={introductory-prerequisite},
  notation={$\HLPopulationSD{\sigma}$},
  examples={},
  source={book/grade10/term_2/1-practice_activities},
  brief={Spread of individual values, in original units.},
  definition={$\sigma_X$ or $\sigma$ is the nonnegative square root of the population or model variance. It describes the spread of individual values in the original units; it is not the standard error of a mean.}
}

\HighlightDeclareContext{g10-t2-intro-size}{
  concept={design-count},
  grade={10},
  term={2},
  criteria={},
  variant={introductory-prerequisite},
  notation={$\HLObservationCount{n}$},
  examples={},
  source={book/grade10/term_2/1-practice_activities},
  brief={Fixed trial count or sample size; a positive integer.}
}

\HighlightDeclareContext{g10-t2-intro-assume}{
  concept={assume},
  grade={10},
  term={2},
  criteria={},
  variant={introductory-prerequisite},
  notation={Bernoulli, binomial, geometric, negative binomial, Poisson},
  examples={},
  source={book/grade10/term_2/1-practice_activities},
  brief={Conditions to check before using the selected distribution.}
}

\HighlightDeclareContext{g10-t2-intro-target}{
  concept={target},
  grade={10},
  term={2},
  criteria={},
  variant={introductory-prerequisite},
  notation={$\HLModelControl{r}$},
  examples={},
  source={book/grade10/term_2/1-practice_activities},
  brief={Fixed number of successes required before stopping.}
}

\HighlightDeclareContext{g10-t2-default-outcome}{
  concept={outcome-space},
  grade={10},
  term={2},
  criteria={},
  variant={compatibility-default},
  notation={$\GTenHLOutcome{\omega}$, $\GTenHLOutcome{\Omega}$},
  examples={},
  source={book/grade10/term_2/1-practice_activities},
  brief={One complete result; all possible complete results.},
  definition={One complete result; all possible complete results.}
}

\HighlightDeclareContext{g10-t2-default-event}{
  concept={event},
  grade={10},
  term={2},
  criteria={},
  variant={compatibility-default},
  notation={$\HLMathEvent{A}$, $\HLMathEvent{E}$, $\HLMathEvent{E^c}$; $\{\GTenHLValue{X}\leq \GTenHLValue{a}\}$},
  examples={},
  source={book/grade10/term_2/1-practice_activities},
  brief={Outcomes meeting a condition; those outside it.},
  definition={Outcomes meeting a condition; those outside it.}
}

\HighlightDeclareContext{g10-t2-default-variable}{
  concept={random-variable},
  grade={10},
  term={2},
  criteria={},
  variant={compatibility-default},
  notation={$\GTenHLValue{X}$, $\GTenHLValue{x}$, $\GTenHLValue{T}$, $\GTenHLValue{t}$; $\GTenHLValue{X}(\GTenHLOutcome{\omega})$},
  examples={},
  source={book/grade10/term_2/1-practice_activities},
  name={\GTenHLValue{Primary} \GTenHLValue{variable} \GTenHLValue{and} \GTenHLValue{value}},
  brief={Numerical rule and possible value; state the units.},
  definition={Numerical rule and possible value; state the units.}
}

\HighlightDeclareContext{g10-t2-default-support}{
  concept={discrete-support},
  grade={10},
  term={2},
  criteria={},
  variant={compatibility-default},
  notation={$\GTenHLValue{S_X}=\{\GTenHLValue{x}:\HLProbability{P}(\GTenHLValue{X}=\GTenHLValue{x})>0\}$},
  examples={},
  source={book/grade10/term_2/1-practice_activities},
  brief={Distinct values with positive probability, in the variable's units.},
  definition={Distinct values with positive probability, in the variable's units.}
}

\HighlightDeclareContext{g10-t2-default-second}{
  concept={second},
  grade={10},
  term={2},
  criteria={},
  variant={compatibility-default},
  notation={$\GTenHLSecond{Y}$, $\GTenHLSecond{y}$; $\GTenHLSecond{Y}=f(\GTenHLValue{X})$; $\GTenHLSecond{Y_i}=f_i(\GTenHLValue{X})$},
  examples={},
  source={book/grade10/term_2/1-practice_activities},
  brief={Another numerical description, with its own units.},
  definition={Another numerical description, with its own units.}
}

\HighlightDeclareContext{g10-t2-default-third}{
  concept={third},
  grade={10},
  term={2},
  criteria={},
  variant={compatibility-default},
  notation={$\GTenHLThird{Z}$, $\GTenHLThird{z}$, $\GTenHLThird{C}$},
  examples={},
  source={book/grade10/term_2/1-practice_activities},
  brief={A third numerical description; an indicator is unitless.},
  definition={A third numerical description; an indicator is unitless.}
}

\HighlightDeclareContext{g10-t2-default-mass}{
  concept={event-probability},
  grade={10},
  term={2},
  criteria={},
  variant={compatibility-default},
  notation={$\HLProbability{P(A)}$, $\HLProbability{P(\omega)}$, $\HLProbability{p_X(x)}$, $\HLProbability{m_X(x)}$},
  examples={},
  source={book/grade10/term_2/1-practice_activities},
  name={\HLProbability{Probability} \HLProbability{and} \HLProbability{mass}},
  brief={Unitless chance of an event or one supported value.},
  definition={Unitless chance of an event or one supported value.}
}

\HighlightDeclareContext{g10-t2-default-unknown}{
  concept={unknown-constant},
  grade={10},
  term={2},
  criteria={},
  variant={compatibility-default},
  notation={$\GTenHLParameter{k}$, $\GTenHLParameter{q}$, $\GTenHLParameter{c}$, $\GTenHLParameter{a}$, $\GTenHLParameter{w}$},
  examples={},
  source={book/grade10/term_2/1-practice_activities},
  name={\GTenHLParameter{Unknown} \GTenHLParameter{quantity}},
  brief={A model quantity to determine, with a context-defined domain.},
  definition={A model quantity to determine, with a context-defined domain.}
}

\HighlightDeclareContext{g10-t2-default-ops}{
  concept={mathematical-operations},
  grade={10},
  term={2},
  criteria={},
  variant={compatibility-default},
  notation={$\sum$, $\binom nk$, $n!$, $e$, $\sim$},
  examples={},
  source={book/grade10/term_2/1-practice_activities},
  brief={Sum, count selections, take factorials, use the exponential constant, or name a model.},
  definition={Sum, count selections, take factorials, use the exponential constant, or name a model.}
}

\HighlightDeclareContext{g10-t2-default-assume}{
  concept={assume},
  grade={10},
  term={2},
  criteria={},
  variant={compatibility-default},
  notation={Independence, constant $p$, fixed observation rule},
  examples={},
  source={book/grade10/term_2/1-practice_activities},
  brief={Conditions to check before using the selected distribution.},
  definition={Conditions to check before using the selected distribution.}
}

\HighlightDeclareContext{g10-t2-default-cum}{
  concept={cumulative-probability},
  grade={10},
  term={2},
  criteria={},
  variant={compatibility-default},
  notation={$\HLCumulativeProbability{F_X(a)}$, $\HLCumulativeProbability{F_X(a^-)}$; $\HLCumulativeProbability{P}(\GTenHLValue{X}\leq \GTenHLValue{a})$},
  examples={},
  source={book/grade10/term_2/1-practice_activities},
  name={\HLCumulativeProbability{Cumulative} \HLCumulativeProbability{probability}},
  brief={Chance up to a cutoff; the left limit excludes the cutoff.},
  definition={Chance up to a cutoff; the left limit excludes the cutoff.}
}

\HighlightDeclareContext{g10-t2-default-cutoff}{
  concept={cutoff},
  grade={10},
  term={2},
  criteria={},
  variant={compatibility-default},
  notation={$\GTenHLValue{a}$, $\GTenHLValue{b}$, $\GTenHLValue{m}$, $\GTenHLValue{\lfloor x\rfloor}$},
  examples={},
  source={book/grade10/term_2/1-practice_activities},
  brief={A boundary in the variable's units; specify included endpoints.},
  definition={A boundary in the variable's units; specify included endpoints.}
}

\HighlightDeclareContext{g10-t2-default-aux}{
  concept={aux},
  grade={10},
  term={2},
  criteria={},
  variant={compatibility-default},
  notation={$K,a,b,c$; $a_m,b_m,c_m$; $N_k,b_j,R_m$},
  examples={},
  source={book/grade10/term_2/1-practice_activities},
  brief={Problem-defined counts, recurrence states, balances or survival quantities.},
  definition={Problem-defined counts, recurrence states, balances or survival quantities.}
}

\HighlightDeclareContext{g10-t2-default-p}{
  concept={success-probability},
  grade={10},
  term={2},
  criteria={},
  variant={compatibility-default},
  notation={$\HLSuccessProbability{p}=\HLSuccessProbability{P}(\GTenHLValue{X}=\GTenHLValue{1})$},
  examples={},
  source={book/grade10/term_2/1-practice_activities},
  name={\HLSuccessProbability{Success} \HLSuccessProbability{probability}},
  brief={Chance of the specified state labelled 1, between 0 and 1.},
  definition={Chance of the specified state labelled 1, between 0 and 1.}
}

\HighlightDeclareContext{g10-t2-default-q}{
  concept={complement-probability},
  grade={10},
  term={2},
  criteria={},
  variant={compatibility-default},
  notation={$\HLComplementProbability{1-p}$},
  examples={},
  source={book/grade10/term_2/1-practice_activities},
  name={\HLComplementProbability{Complement} \HLComplementProbability{probability}},
  brief={Chance of the complementary state labelled 0.},
  definition={Chance of the complementary state labelled 0.}
}

\HighlightDeclareContext{g10-t2-default-expected}{
  concept={expected-value},
  grade={10},
  term={2},
  criteria={},
  variant={compatibility-default},
  notation={$\HLExpectedValue{E}(\GTenHLValue{X})$, $\HLExpectedValue{\mathbb E}[\GTenHLSecond{Y_i}]$},
  examples={},
  source={book/grade10/term_2/1-practice_activities},
  name={\HLExpectedValue{Expected} \HLExpectedValue{value}},
  brief={Probability-weighted average, in original units.},
  definition={Probability-weighted average, in original units.}
}

\HighlightDeclareContext{g10-t2-default-meanpop}{
  concept={population-mean},
  grade={10},
  term={2},
  criteria={},
  variant={compatibility-default},
  notation={$\HLPopulationMean{\mu}$},
  examples={},
  source={book/grade10/term_2/1-practice_activities},
  name={\HLPopulationMean{Population} \HLPopulationMean{or} \HLPopulationMean{model} \HLPopulationMean{mean}},
  brief={Mean of the population or specified probability model, in original units.},
  definition={Mean of the population or specified probability model, in original units.}
}

\HighlightDeclareContext{g10-t2-default-break}{
  concept={break},
  grade={10},
  term={2},
  criteria={},
  variant={compatibility-default},
  notation={$\HLSuccessProbability{p_i^*}$},
  examples={},
  source={book/grade10/term_2/1-practice_activities},
  brief={Probability at which the specified expected net payoff is zero.},
  definition={Probability at which the specified expected net payoff is zero.}
}

\HighlightDeclareContext{g10-t2-default-transform}{
  concept={transform},
  grade={10},
  term={2},
  criteria={},
  variant={compatibility-default},
  notation={$a,b$; $f_i(0),f_i(1)$},
  examples={},
  source={book/grade10/term_2/1-practice_activities},
  brief={Scale, fixed addition or state payoffs, with required units.},
  definition={Scale, fixed addition or state payoffs, with required units.}
}

\HighlightDeclareContext{g10-t2-default-size}{
  concept={design-count},
  grade={10},
  term={2},
  criteria={},
  variant={compatibility-default},
  notation={$\HLObservationCount{n}$},
  examples={},
  source={book/grade10/term_2/1-practice_activities},
  brief={Fixed trial count or sample size; a positive integer.},
  definition={Fixed trial count or sample size; a positive integer.}
}

\HighlightDeclareContext{g10-t2-default-wait}{
  concept={wait},
  grade={10},
  term={2},
  criteria={},
  variant={compatibility-default},
  notation={$\GTenHLValue{T}$, $\GTenHLValue{t}$; $\GTenHLValue{F}=\GTenHLValue{T}-1$ or $\GTenHLValue{F}=\GTenHLValue{T}-\HLModelControl{r}$},
  examples={},
  source={book/grade10/term_2/1-practice_activities},
  brief={Total trials include the final success; failures exclude successes.},
  definition={Total trials include the final success; failures exclude successes.}
}

\HighlightDeclareContext{g10-t2-default-target}{
  concept={target},
  grade={10},
  term={2},
  criteria={},
  variant={compatibility-default},
  notation={$\HLModelControl{r}$},
  examples={},
  source={book/grade10/term_2/1-practice_activities},
  brief={Fixed number of successes required before stopping.},
  definition={Fixed number of successes required before stopping.}
}

\HighlightDeclareContext{g10-t2-default-rate}{
  concept={event-rate-per-exposure},
  grade={10},
  term={2},
  criteria={},
  variant={compatibility-default},
  notation={$\mathrm{rate}$, $\mathrm{exposure}$},
  examples={},
  source={book/grade10/term_2/1-practice_activities},
  brief={Events per exposure unit and length of the observation window.},
  definition={Events per exposure unit and length of the observation window.}
}

\HighlightDeclareContext{g10-t2-default-lambda}{
  concept={poisson-expected-count},
  grade={10},
  term={2},
  criteria={},
  variant={compatibility-default},
  notation={$\HLPoissonMean{\lambda}$},
  examples={},
  source={book/grade10/term_2/1-practice_activities},
  name={\HLPoissonMean{Poisson} \HLPoissonMean{interval} \HLPoissonMean{mean}},
  brief={Expected count for the selected exposure.},
  definition={Expected count for the selected exposure.}
}

\HighlightDeclareContext{g10-t2-default-secondmoment}{
  concept={secondmoment},
  grade={10},
  term={2},
  criteria={},
  variant={compatibility-default},
  notation={$\HLExpectedValue{E}(\GTenHLValue{X^2})$, $[\HLExpectedValue{E}(\GTenHLValue{X})]^2$},
  examples={},
  source={book/grade10/term_2/1-practice_activities},
  brief={Average of squares; square of the average. Both use squared units.},
  definition={Average of squares; square of the average. Both use squared units.}
}

\HighlightDeclareContext{g10-t2-default-variance}{
  concept={population-variance},
  grade={10},
  term={2},
  criteria={},
  variant={compatibility-default},
  notation={$\HLPopulationVariance{\operatorname{Var}}(\GTenHLValue{X})$, $\HLPopulationVariance{\sigma_X^2}$},
  examples={},
  source={book/grade10/term_2/1-practice_activities},
  name={\HLPopulationVariance{Variance}},
  brief={Mean squared deviation from the mean; squared units.},
  definition={Mean squared deviation from the mean; squared units.}
}

\HighlightDeclareContext{g10-t2-default-sd}{
  concept={population-standard-deviation},
  grade={10},
  term={2},
  criteria={},
  variant={compatibility-default},
  notation={$\HLPopulationSD{\sigma_X}$, $\HLPopulationSD{\sigma}$},
  examples={},
  source={book/grade10/term_2/1-practice_activities},
  name={\HLPopulationSD{Population} \HLPopulationSD{or} \HLPopulationSD{model} \HLPopulationSD{standard} \HLPopulationSD{deviation}},
  brief={Spread of individual values, in original units.},
  definition={Spread of individual values, in original units.}
}

\HighlightDeclareContext{g10-t2-legacy-cell-1}{
  concept={outcome-space},
  grade={10},
  term={2},
  criteria={1},
  variant={legacy-cell},
  notation={\GTenHLOutcome{Outcome}},
  examples={},
  source={highlights_compat_ttwo.tex:G10-T2-C1},
  display={\GTenHLOutcome{Outcome}},
  meaning={A possible result of the experiment.}
}

\HighlightDeclareContext{g10-t2-legacy-cell-2}{
  concept={event-probability},
  grade={10},
  term={2},
  criteria={1},
  variant={legacy-cell},
  notation={\GTenHLProbability{Probability}},
  examples={},
  source={highlights_compat_ttwo.tex:G10-T2-C1},
  display={\GTenHLProbability{Probability}},
  meaning={The chance of an outcome or event, expressed without units.}
}

\HighlightDeclareContext{g10-t2-legacy-cell-3}{
  concept={random-variable},
  grade={10},
  term={2},
  criteria={1},
  variant={legacy-cell},
  notation={\GTenHLValue{First variable}},
  examples={},
  source={highlights_compat_ttwo.tex:G10-T2-C1},
  display={\GTenHLValue{First variable}},
  meaning={The first variable, its numerical values, and any units.}
}

\HighlightDeclareContext{g10-t2-legacy-cell-4}{
  concept={second},
  grade={10},
  term={2},
  criteria={1},
  variant={legacy-cell},
  notation={\GTenHLSecond{Second variable}},
  examples={},
  source={highlights_compat_ttwo.tex:G10-T2-C1},
  display={\GTenHLSecond{Second variable}},
  meaning={The second variable, its numerical values, and any units.}
}

\HighlightDeclareContext{g10-t2-legacy-cell-5}{
  concept={third},
  grade={10},
  term={2},
  criteria={1},
  variant={legacy-cell},
  notation={\GTenHLThird{Third variable}},
  examples={},
  source={highlights_compat_ttwo.tex:G10-T2-C1},
  display={\GTenHLThird{Third variable}},
  meaning={The third variable, its numerical values, and any units.}
}

\HighlightDeclareContext{g10-t2-legacy-cell-6}{
  concept={unknown-constant},
  grade={10},
  term={2},
  criteria={1},
  variant={legacy-cell},
  notation={\GTenHLParameter{Unknown parameter}},
  examples={},
  source={highlights_compat_ttwo.tex:G10-T2-C1},
  display={\GTenHLParameter{Unknown parameter}},
  meaning={A quantity to determine from the information given.}
}

\HighlightDeclareContext{g10-t2-legacy-cell-7}{
  concept={random-variable},
  grade={10},
  term={2},
  criteria={8},
  variant={legacy-cell},
  notation={$\GTenHLValue{X},\ \GTenHLValue{x}$},
  examples={},
  source={highlights_compat_ttwo.tex:G10-T2-C8},
  display={},
  meaning={Quantity and supported values; retain their units.},
  name={\GTenHLValue{Primary variable}}
}

\HighlightDeclareContext{g10-t2-legacy-cell-8}{
  concept={second},
  grade={10},
  term={2},
  criteria={8},
  variant={legacy-cell},
  notation={$\GTenHLSecond{Y}$},
  examples={},
  source={highlights_compat_ttwo.tex:G10-T2-C8},
  display={},
  meaning={Quantity after scaling and a fixed addition.},
  name={\GTenHLSecond{Transformed variable}}
}

\HighlightDeclareContext{g10-t2-legacy-cell-9}{
  concept={expected-value},
  grade={10},
  term={2},
  criteria={8},
  variant={legacy-cell},
  notation={$\HLExpectedValue{E}(\GTenHLValue{X})$, $\HLExpectedValue{E}(\GTenHLSecond{Y})$},
  examples={},
  source={highlights_compat_ttwo.tex:G10-T2-C8},
  display={},
  meaning={Probability-weighted average; original units.},
  name={\HLExpectedValue{Expected value}}
}

\HighlightDeclareContext{g10-t2-legacy-cell-10}{
  concept={secondmoment},
  grade={10},
  term={2},
  criteria={8},
  variant={legacy-cell},
  notation={$\HLExpectedValue{E}(\GTenHLValue{X^2})$},
  examples={},
  source={highlights_compat_ttwo.tex:G10-T2-C8},
  display={},
  meaning={Average after squaring; squared units.},
  name={\HLExpectedValue{Expected value} of squared values}
}


% Opt-in notation contexts under existing concepts; concise unions stay grouped.
\HighlightDeclareContext{g10-t2-c1-notation-outcome}{
  concept={outcome-space},
  role={outcome},
  grade={10},
  term={2},
  criteria={1},
  variant={notation},
  notation={$\GTenHLOutcome{\omega}$},
  name={\GTenHLOutcome{Outcome}},
  brief={One complete result of an experiment; an element of the sample space $\Omega$.},
  definition={One complete result of an experiment; an element of the sample space $\Omega$.},
  examples={},
  source={book/grade10/term_2/1-practice_activities/G10_T2_L1_C1_practice_activity.tex}
}

\HighlightDeclareContext{g10-t2-c1-notation-variable}{
  concept={random-variable},
  role={quantity.primary},
  grade={10},
  term={2},
  criteria={1},
  variant={notation},
  notation={$\GTenHLValue{X}$},
  name={\GTenHLValue{Random variable}},
  brief={$X$ is a numerical rule assigning a value to each complete outcome; the model specifies its meaning and units.},
  definition={$X$ is a numerical rule assigning a value to each complete outcome; the model specifies its meaning and units.},
  examples={},
  source={book/grade10/term_2/1-practice_activities/G10_T2_L1_C1_practice_activity.tex}
}

\HighlightDeclareContext{g10-t2-c1-notation-value}{
  concept={random-variable},
  role={quantity.primary},
  grade={10},
  term={2},
  criteria={1},
  variant={notation},
  notation={$\GTenHLValue{x}$},
  name={\GTenHLValue{Variable value}},
  brief={$x$ denotes a numerical value in the units of $X$. In a support table it is a possible value; a PMF or CDF may be evaluated at any real argument, including one outside the support.},
  definition={$x$ denotes a numerical value in the units of $X$. In a support table it is a possible value; a PMF or CDF may be evaluated at any real argument, including one outside the support.},
  examples={},
  source={book/grade10/term_2/1-practice_activities/G10_T2_L1_C1_practice_activity.tex}
}

\HighlightDeclareContext{g10-t2-c2-notation-outcome}{
  concept={outcome-space},
  role={outcome},
  grade={10},
  term={2},
  criteria={2},
  variant={notation},
  notation={$\GTenHLOutcome{\omega}$},
  name={\GTenHLOutcome{Outcome}},
  brief={One complete result of an experiment; an element of the sample space $\Omega$.},
  definition={One complete result of an experiment; an element of the sample space $\Omega$.},
  examples={},
  source={book/grade10/term_2/1-practice_activities/G10_T2_L1_C2_practice_activity.tex}
}

\HighlightDeclareContext{g10-t2-c2-notation-variable}{
  concept={random-variable},
  role={quantity.primary},
  grade={10},
  term={2},
  criteria={2},
  variant={notation},
  notation={$\GTenHLValue{X}$},
  name={\GTenHLValue{Random variable}},
  brief={$X$ is a numerical rule assigning a value to each complete outcome; the model specifies its meaning and units.},
  definition={$X$ is a numerical rule assigning a value to each complete outcome; the model specifies its meaning and units.},
  examples={},
  source={book/grade10/term_2/1-practice_activities/G10_T2_L1_C2_practice_activity.tex}
}

\HighlightDeclareContext{g10-t2-c2-notation-value}{
  concept={random-variable},
  role={quantity.primary},
  grade={10},
  term={2},
  criteria={2},
  variant={notation},
  notation={$\GTenHLValue{x}$},
  name={\GTenHLValue{Variable value}},
  brief={$x$ denotes a numerical value in the units of $X$. In a support table it is a possible value; a PMF or CDF may be evaluated at any real argument, including one outside the support.},
  definition={$x$ denotes a numerical value in the units of $X$. In a support table it is a possible value; a PMF or CDF may be evaluated at any real argument, including one outside the support.},
  examples={},
  source={book/grade10/term_2/1-practice_activities/G10_T2_L1_C2_practice_activity.tex}
}

\HighlightDeclareContext{g10-t2-c1-notation-space}{
  concept={outcome-space},
  role={outcome},
  grade={10},
  term={2},
  criteria={1},
  variant={notation},
  notation={$\GTenHLOutcome{\Omega}$},
  name={\GTenHLOutcome{Sample space}},
  brief={The set of all possible complete outcomes of the experiment.},
  definition={The set of all possible complete outcomes of the experiment.},
  examples={},
  source={book/grade10/term_2/1-practice_activities/G10_T2_L1_C1_practice_activity.tex}
}

\HighlightDeclareContext{g10-t2-c2-notation-inclusive}{
  concept={cumulative-probability},
  role={cumulative},
  grade={10},
  term={2},
  criteria={2},
  variant={notation},
  notation={$\HLCumulativeProbability{F_X(a)}$},
  name={\HLCumulativeProbability{Inclusive cumulative height}},
  brief={$F_X(a)=P(X\leq a)$ includes all supported values no greater than the real cutoff $a$. The output is unitless; the cutoff has the units of $X$ and need not be supported.},
  definition={$F_X(a)=P(X\leq a)$ includes all supported values no greater than the real cutoff $a$. The output is unitless; the cutoff has the units of $X$ and need not be supported.},
  examples={},
  source={book/grade10/term_2/1-practice_activities/G10_T2_L1_C2_practice_activity.tex}
}

\HighlightDeclareContext{g10-t2-c2-notation-left-limit}{
  concept={cumulative-probability},
  role={cumulative},
  grade={10},
  term={2},
  criteria={2},
  variant={notation},
  notation={$\HLCumulativeProbability{F_X(a^-)}$},
  name={\HLCumulativeProbability{Left-limit height}},
  brief={$F_X(a^-)=P(X<a)$ is the cumulative probability strictly before $a$, excluding equality; it differs from the inclusive value $F_X(a)$.},
  definition={$F_X(a^-)=P(X<a)$ is the cumulative probability strictly before $a$, excluding equality; it differs from the inclusive value $F_X(a)$.},
  examples={},
  source={book/grade10/term_2/1-practice_activities/G10_T2_L1_C2_practice_activity.tex}
}

\HighlightDeclareContext{g10-t2-c2-notation-jump}{
  concept={event-probability},
  role={probability},
  grade={10},
  term={2},
  criteria={2},
  variant={notation},
  notation={$\HLCumulativeProbability{F_X(a)}-\HLCumulativeProbability{F_X(a^-)}$},
  name={\HLProbability{Jump size}},
  brief={For a discrete variable, $\HLCumulativeProbability{F_X(a)}-\HLCumulativeProbability{F_X(a^-)}=\HLProbability{p_X(a)}=P(X=a)$. A jump is one unitless point mass, rather than a cumulative height; it is zero at an unsupported cutoff.},
  definition={For a discrete variable, $\HLCumulativeProbability{F_X(a)}-\HLCumulativeProbability{F_X(a^-)}=\HLProbability{p_X(a)}=P(X=a)$. A jump is one unitless point mass, rather than a cumulative height; it is zero at an unsupported cutoff.},
  examples={},
  source={book/grade10/term_2/1-practice_activities/G10_T2_L1_C2_practice_activity.tex}
}

\HighlightDeclareContext{g10-t2-c3-notation-function}{
  concept={transform},
  role={transformation.function.c3},
  grade={10},
  term={2},
  criteria={3},
  variant={notation},
  notation={$\HLGtenCThreeFunction{f}$},
  name={\HLGtenCThreeFunction{Transformation function}},
  brief={A rule assigning an output value to each input state, with $Y=\HLGtenCThreeFunction{f}(X)$. Outputs have the units of $Y$. The rule is distinct from a PMF or CDF.},
  definition={A rule assigning an output value to each input state, with $Y=\HLGtenCThreeFunction{f}(X)$. Outputs have the units of $Y$. The rule is distinct from a PMF or CDF.},
  examples={},
  source={book/grade10/term_2/1-practice_activities/G10_T2_L1_C3_practice_activity.tex}
}

\HighlightDeclareContext{g10-t2-c3-notation-indexed-function}{
  concept={transform},
  role={transformation.function.c3},
  grade={10},
  term={2},
  criteria={3},
  variant={notation},
  notation={$\HLGtenCThreeFunction{f_i}$},
  name={\HLGtenCThreeFunction{Indexed transformation function}},
  brief={$\HLGtenCThreeFunction{f_i}$ is the rule for alternative $i$, with $Y_i=\HLGtenCThreeFunction{f_i}(X)$. The index identifies the rule, rather than a function evaluation or probability.},
  definition={$\HLGtenCThreeFunction{f_i}$ is the rule for alternative $i$, with $Y_i=\HLGtenCThreeFunction{f_i}(X)$. The index identifies the rule, rather than a function evaluation or probability.},
  examples={},
  source={book/grade10/term_2/1-practice_activities/G10_T2_L1_C3_practice_activity.tex}
}

\HighlightDeclareContext{g10-t2-c3-notation-payoff-0}{
  concept={transform},
  role={quantity.secondary},
  grade={10},
  term={2},
  criteria={3},
  variant={notation},
  notation={$\GTenHLSecond{f_i(0)}$},
  name={\GTenHLSecond{State-0 payoff}},
  brief={The output value assigned by $f_i$ at input state $0$, in the units of $Y_i$. It is not a probability; its sign depends on the context.},
  definition={The output value assigned by $f_i$ at input state $0$, in the units of $Y_i$. It is not a probability; its sign depends on the context.},
  examples={},
  source={book/grade10/term_2/1-practice_activities/G10_T2_L1_C3_practice_activity.tex}
}

\HighlightDeclareContext{g10-t2-c3-notation-payoff-1}{
  concept={transform},
  role={quantity.secondary},
  grade={10},
  term={2},
  criteria={3},
  variant={notation},
  notation={$\GTenHLSecond{f_i(1)}$},
  name={\GTenHLSecond{State-1 payoff}},
  brief={The output value assigned by $f_i$ at input state $1$, in the units of $Y_i$. It is not a probability; its sign depends on the context.},
  definition={The output value assigned by $f_i$ at input state $1$, in the units of $Y_i$. It is not a probability; its sign depends on the context.},
  examples={},
  source={book/grade10/term_2/1-practice_activities/G10_T2_L1_C3_practice_activity.tex}
}

\HighlightDeclareContext{g10-t2-c3-notation-indicator-mass}{
  concept={event-probability},
  role={probability},
  grade={10},
  term={2},
  criteria={3},
  variant={notation},
  notation={$\HLProbability{m_X(x)}$},
  name={\HLProbability{Indicator PMF}},
  brief={For $X\sim\operatorname{Bernoulli}(p)$, $m_X(x)=P(X=x)$ is $1-p$ at 0, $p$ at 1 and zero elsewhere, including at $p=0$ or $p=1$. The output is unitless.},
  definition={For $X\sim\operatorname{Bernoulli}(p)$, $m_X(x)=P(X=x)$ is $1-p$ at 0, $p$ at 1 and zero elsewhere, including at $p=0$ or $p=1$. The output is unitless.},
  examples={},
  source={book/grade10/term_2/1-practice_activities/G10_T2_L1_C3_practice_activity.tex}
}

\HighlightDeclareContext{g10-t2-c3-notation-transformed-mass}{
  concept={event-probability},
  role={probability},
  grade={10},
  term={2},
  criteria={3},
  variant={notation},
  notation={$\HLProbability{m_{Y_i}(y)}$},
  name={\HLProbability{Transformed PMF}},
  brief={$m_{Y_i}(y)=P(Y_i=y)$ assigns unitless probability to a value in the units of $Y_i$. For $Y_i=f_i(X)$, sum all state masses mapping to $y$; the mass is zero if no positive-mass state maps there. This is distinct from $f_i$ or a density.},
  definition={$m_{Y_i}(y)=P(Y_i=y)$ assigns unitless probability to a value in the units of $Y_i$. For $Y_i=f_i(X)$, sum all state masses mapping to $y$; the mass is zero if no positive-mass state maps there. This is distinct from $f_i$ or a density.},
  examples={},
  source={book/grade10/term_2/1-practice_activities/G10_T2_L1_C3_practice_activity.tex}
}

\HighlightDeclareContext{g10-t2-c3-notation-indicator-cdf}{
  concept={cumulative-probability},
  role={cumulative},
  grade={10},
  term={2},
  criteria={3},
  variant={notation},
  notation={$\HLCumulativeProbability{F_X(t)}$},
  name={\HLCumulativeProbability{Indicator CDF}},
  brief={For a Bernoulli indicator, $F_X(t)=P(X\leq t)$ equals 0 for $t<0$, $1-p$ for $0\leq t<1$, and 1 for $t\geq1$. It is right-continuous; threshold and output are unitless.},
  definition={For a Bernoulli indicator, $F_X(t)=P(X\leq t)$ equals 0 for $t<0$, $1-p$ for $0\leq t<1$, and 1 for $t\geq1$. It is right-continuous; threshold and output are unitless.},
  examples={},
  source={book/grade10/term_2/1-practice_activities/G10_T2_L1_C3_practice_activity.tex}
}

\HighlightDeclareContext{g10-t2-c3-notation-transformed-cdf}{
  concept={cumulative-probability},
  role={cumulative},
  grade={10},
  term={2},
  criteria={3},
  variant={notation},
  notation={$\HLCumulativeProbability{F_{Y_i}(t)}$},
  name={\HLCumulativeProbability{Transformed CDF}},
  brief={$F_{Y_i}(t)=P(Y_i\leq t)$ accumulates masses in numerical value order, including the cutoff. It is right-continuous; $t$ has the units of $Y_i$, while the output is unitless. State labels do not determine payoff order.},
  definition={$F_{Y_i}(t)=P(Y_i\leq t)$ accumulates masses in numerical value order, including the cutoff. It is right-continuous; $t$ has the units of $Y_i$, while the output is unitless. State labels do not determine payoff order.},
  examples={},
  source={book/grade10/term_2/1-practice_activities/G10_T2_L1_C3_practice_activity.tex}
}

\HighlightDeclareContext{g10-t2-c3-notation-expectation-operator}{
  concept={expected-value},
  role={expectation},
  grade={10},
  term={2},
  criteria={3},
  variant={notation},
  notation={$\HLExpectedValue{\mathbb E}$},
  name={\HLExpectedValue{Expectation operator}},
  brief={Forms a probability-weighted average of a random variable, when it exists. The result retains the units of the variable. The operator differs from an event or indicator named $E$.},
  definition={Forms a probability-weighted average of a random variable, when it exists. The result retains the units of the variable. The operator differs from an event or indicator named $E$.},
  examples={},
  source={book/grade10/term_2/1-practice_activities/G10_T2_L1_C3_practice_activity.tex}
}

\HighlightDeclareReference{G10-T2-C1-notation}{
  status={available},
  members={g10-t2-c1-notation-outcome,g10-t2-c1-notation-variable,g10-t2-c1-notation-value,g10-t2-c1-notation-space}
}

\HighlightDeclareReference{G10-T2-C2-notation}{
  status={available},
  members={g10-t2-c2-notation-outcome,g10-t2-c2-notation-variable,g10-t2-c2-notation-value,g10-t2-c2-notation-inclusive,g10-t2-c2-notation-left-limit,g10-t2-c2-notation-jump}
}

\HighlightDeclareReference{G10-T2-C3-notation}{
  status={available},
  members={g10-t2-c3-notation-function,g10-t2-c3-notation-indexed-function,g10-t2-c3-notation-payoff-0,g10-t2-c3-notation-payoff-1,g10-t2-c3-notation-indicator-mass,g10-t2-c3-notation-transformed-mass,g10-t2-c3-notation-indicator-cdf,g10-t2-c3-notation-transformed-cdf,g10-t2-c3-notation-expectation-operator}
}

\HighlightDeclareReference{G10-T2-C1-C2-C3-C4-C5-C6-C7-C8-C9}{
  status={available},
  members={g10-t2-c1-c2-c3-c4-c5-c6-c7-c8-c9-10-outcome,g10-t2-c1-c2-c3-c4-c5-c6-c7-c8-c9-10-mass,g10-t2-c1-c2-c3-c4-c5-c6-c7-c8-c9-10-variable,g10-t2-c1-c2-c3-c4-c5-c6-c7-c8-c9-10-second,g10-t2-c1-c2-c3-c4-c5-c6-c7-c8-c9-10-third,g10-t2-c1-c2-c3-c4-c5-c6-c7-c8-c9-10-unknown},
  renderer={G10-T2-C1-C2-C3-C4-C5-C6-C7-C8-C9}
}

\HighlightDeclareReference{G10-T2-C8}{
  status={available},
  members={g10-t2-c8-10-variable,g10-t2-c8-10-second,g10-t2-c8-10-expected,g10-t2-c8-10-secondmoment,g10-t2-legacy-cell-7,g10-t2-legacy-cell-8,g10-t2-legacy-cell-9,g10-t2-legacy-cell-10},
  renderer={G10-T2-C8}
}

\HighlightDeclareReference{G10-T2-C1-C2-C3-C4-C5-C6-C7-C8-C9-coordinated}{
  status={available},
  members={g10-t2-c1-coordinated-10-outcome,g10-t2-c1-coordinated-10-event,g10-t2-c1-coordinated-10-variable,g10-t2-c1-coordinated-10-support,g10-t2-c1-coordinated-10-second,g10-t2-c1-coordinated-10-third,g10-t2-c1-coordinated-10-mass,g10-t2-c1-coordinated-10-unknown,g10-t2-c1-coordinated-10-ops,g10-t2-c1-coordinated-10-assume,g10-t2-c2-coordinated-10-outcome,g10-t2-c2-coordinated-10-event,g10-t2-c2-coordinated-10-variable,g10-t2-c2-coordinated-10-support,g10-t2-c2-coordinated-10-mass,g10-t2-c2-coordinated-10-cum,g10-t2-c2-coordinated-10-unknown,g10-t2-c2-coordinated-10-cutoff,g10-t2-c2-coordinated-10-ops,g10-t2-c2-coordinated-10-assume,g10-t2-c2-coordinated-10-aux,g10-t2-c3-coordinated-10-event,g10-t2-c3-coordinated-10-variable,g10-t2-c3-coordinated-10-support,g10-t2-c3-coordinated-10-p,g10-t2-c3-coordinated-10-q,g10-t2-c3-coordinated-10-second,g10-t2-c3-coordinated-10-mass,g10-t2-c3-coordinated-10-cum,g10-t2-c3-coordinated-10-expected,g10-t2-c3-coordinated-10-break,g10-t2-c3-coordinated-10-transform,g10-t2-c4-coordinated-10-variable,g10-t2-c4-coordinated-10-support,g10-t2-c4-coordinated-10-second,g10-t2-c4-coordinated-10-p,g10-t2-c4-coordinated-10-q,g10-t2-c4-coordinated-10-size,g10-t2-c4-coordinated-10-mass,g10-t2-c4-coordinated-10-cum,g10-t2-c4-coordinated-10-cutoff,g10-t2-c4-coordinated-10-ops,g10-t2-c4-coordinated-10-assume,g10-t2-c5-coordinated-10-wait,g10-t2-c5-coordinated-10-support,g10-t2-c5-coordinated-10-p,g10-t2-c5-coordinated-10-q,g10-t2-c5-coordinated-10-second,g10-t2-c5-coordinated-10-mass,g10-t2-c5-coordinated-10-cum,g10-t2-c5-coordinated-10-cutoff,g10-t2-c5-coordinated-10-event,g10-t2-c5-coordinated-10-assume,g10-t2-c6-coordinated-10-wait,g10-t2-c6-coordinated-10-support,g10-t2-c6-coordinated-10-p,g10-t2-c6-coordinated-10-q,g10-t2-c6-coordinated-10-target,g10-t2-c6-coordinated-10-second,g10-t2-c6-coordinated-10-mass,g10-t2-c6-coordinated-10-cum,g10-t2-c6-coordinated-10-size,g10-t2-c6-coordinated-10-cutoff,g10-t2-c6-coordinated-10-ops,g10-t2-c6-coordinated-10-assume,g10-t2-c7-coordinated-10-variable,g10-t2-c7-coordinated-10-support,g10-t2-c7-coordinated-10-mass,g10-t2-c7-coordinated-10-cum,g10-t2-c7-coordinated-10-cutoff,g10-t2-c7-coordinated-10-rate,g10-t2-c7-coordinated-10-lambda,g10-t2-c7-coordinated-10-second,g10-t2-c7-coordinated-10-ops,g10-t2-c7-coordinated-10-assume,g10-t2-c8-coordinated-10-variable,g10-t2-c8-coordinated-10-support,g10-t2-c8-coordinated-10-second,g10-t2-c8-coordinated-10-mass,g10-t2-c8-coordinated-10-expected,g10-t2-c8-coordinated-10-meanpop,g10-t2-c8-coordinated-10-secondmoment,g10-t2-c8-coordinated-10-variance,g10-t2-c8-coordinated-10-sd,g10-t2-c8-coordinated-10-transform,g10-t2-c8-coordinated-10-ops,g10-t2-c9-coordinated-10-event,g10-t2-c9-coordinated-10-variable,g10-t2-c9-coordinated-10-support,g10-t2-c9-coordinated-10-wait,g10-t2-c9-coordinated-10-second,g10-t2-c9-coordinated-10-p,g10-t2-c9-coordinated-10-q,g10-t2-c9-coordinated-10-size,g10-t2-c9-coordinated-10-target,g10-t2-c9-coordinated-10-mass,g10-t2-c9-coordinated-10-cum,g10-t2-c9-coordinated-10-rate,g10-t2-c9-coordinated-10-lambda,g10-t2-c9-coordinated-10-ops,g10-t2-c9-coordinated-10-assume},
  renderer={G10-T2-C1-C2-C3-C4-C5-C6-C7-C8-C9-coordinated}
}

\HighlightDeclareReference{G10-T2-C1-coordinated}{
  status={available},
  members={g10-t2-c1-coordinated-10-outcome,g10-t2-c1-coordinated-10-event,g10-t2-c1-coordinated-10-variable,g10-t2-c1-coordinated-10-support,g10-t2-c1-coordinated-10-second,g10-t2-c1-coordinated-10-third,g10-t2-c1-coordinated-10-mass,g10-t2-c1-coordinated-10-unknown,g10-t2-c1-coordinated-10-ops,g10-t2-c1-coordinated-10-assume},
  renderer={G10-T2-C1-coordinated}
}

\HighlightDeclareReference{G10-T2-C2-coordinated}{
  status={available},
  members={g10-t2-c2-coordinated-10-outcome,g10-t2-c2-coordinated-10-event,g10-t2-c2-coordinated-10-variable,g10-t2-c2-coordinated-10-support,g10-t2-c2-coordinated-10-mass,g10-t2-c2-coordinated-10-cum,g10-t2-c2-coordinated-10-unknown,g10-t2-c2-coordinated-10-cutoff,g10-t2-c2-coordinated-10-ops,g10-t2-c2-coordinated-10-assume,g10-t2-c2-coordinated-10-aux},
  renderer={G10-T2-C2-coordinated}
}

\HighlightDeclareReference{G10-T2-C3-coordinated}{
  status={available},
  members={g10-t2-c3-coordinated-10-event,g10-t2-c3-coordinated-10-variable,g10-t2-c3-coordinated-10-support,g10-t2-c3-coordinated-10-p,g10-t2-c3-coordinated-10-q,g10-t2-c3-coordinated-10-second,g10-t2-c3-coordinated-10-function,g10-t2-c3-coordinated-10-mass,g10-t2-c3-coordinated-10-cum,g10-t2-c3-coordinated-10-expected,g10-t2-c3-coordinated-10-break,g10-t2-c3-coordinated-10-transform},
  renderer={G10-T2-C3-coordinated}
}

\HighlightDeclareReference{G10-T2-C4-coordinated}{
  status={available},
  members={g10-t2-c4-coordinated-10-variable,g10-t2-c4-coordinated-10-support,g10-t2-c4-coordinated-10-second,g10-t2-c4-coordinated-10-p,g10-t2-c4-coordinated-10-q,g10-t2-c4-coordinated-10-size,g10-t2-c4-coordinated-10-mass,g10-t2-c4-coordinated-10-cum,g10-t2-c4-coordinated-10-cutoff,g10-t2-c4-coordinated-10-ops,g10-t2-c4-coordinated-10-assume},
  renderer={G10-T2-C4-coordinated}
}

\HighlightDeclareReference{G10-T2-C5-coordinated}{
  status={available},
  members={g10-t2-c5-coordinated-10-wait,g10-t2-c5-coordinated-10-support,g10-t2-c5-coordinated-10-p,g10-t2-c5-coordinated-10-q,g10-t2-c5-coordinated-10-second,g10-t2-c5-coordinated-10-mass,g10-t2-c5-coordinated-10-cum,g10-t2-c5-coordinated-10-cutoff,g10-t2-c5-coordinated-10-event,g10-t2-c5-coordinated-10-assume},
  renderer={G10-T2-C5-coordinated}
}

\HighlightDeclareReference{G10-T2-C6-coordinated}{
  status={available},
  members={g10-t2-c6-coordinated-10-wait,g10-t2-c6-coordinated-10-support,g10-t2-c6-coordinated-10-p,g10-t2-c6-coordinated-10-q,g10-t2-c6-coordinated-10-target,g10-t2-c6-coordinated-10-second,g10-t2-c6-coordinated-10-mass,g10-t2-c6-coordinated-10-cum,g10-t2-c6-coordinated-10-size,g10-t2-c6-coordinated-10-cutoff,g10-t2-c6-coordinated-10-ops,g10-t2-c6-coordinated-10-assume},
  renderer={G10-T2-C6-coordinated}
}

\HighlightDeclareReference{G10-T2-C7-coordinated}{
  status={available},
  members={g10-t2-c7-coordinated-10-variable,g10-t2-c7-coordinated-10-support,g10-t2-c7-coordinated-10-mass,g10-t2-c7-coordinated-10-cum,g10-t2-c7-coordinated-10-cutoff,g10-t2-c7-coordinated-10-rate,g10-t2-c7-coordinated-10-lambda,g10-t2-c7-coordinated-10-second,g10-t2-c7-coordinated-10-ops,g10-t2-c7-coordinated-10-assume},
  renderer={G10-T2-C7-coordinated}
}

\HighlightDeclareReference{G10-T2-C8-coordinated}{
  status={available},
  members={g10-t2-c8-coordinated-10-variable,g10-t2-c8-coordinated-10-support,g10-t2-c8-coordinated-10-second,g10-t2-c8-coordinated-10-mass,g10-t2-c8-coordinated-10-expected,g10-t2-c8-coordinated-10-meanpop,g10-t2-c8-coordinated-10-secondmoment,g10-t2-c8-coordinated-10-variance,g10-t2-c8-coordinated-10-sd,g10-t2-c8-coordinated-10-transform,g10-t2-c8-coordinated-10-ops},
  renderer={G10-T2-C8-coordinated}
}

\HighlightDeclareReference{G10-T2-C9-coordinated}{
  status={available},
  members={g10-t2-c9-coordinated-10-event,g10-t2-c9-coordinated-10-variable,g10-t2-c9-coordinated-10-support,g10-t2-c9-coordinated-10-wait,g10-t2-c9-coordinated-10-second,g10-t2-c9-coordinated-10-p,g10-t2-c9-coordinated-10-q,g10-t2-c9-coordinated-10-size,g10-t2-c9-coordinated-10-target,g10-t2-c9-coordinated-10-mass,g10-t2-c9-coordinated-10-cum,g10-t2-c9-coordinated-10-rate,g10-t2-c9-coordinated-10-lambda,g10-t2-c9-coordinated-10-ops,g10-t2-c9-coordinated-10-assume},
  renderer={G10-T2-C9-coordinated}
}

\HighlightDeclareReference{G10-T2-C1}{
  status={available},
  members={g10-t2-c1-10-outcome,g10-t2-c1-10-mass,g10-t2-c1-10-variable,g10-t2-c1-10-second,g10-t2-c1-10-third,g10-t2-c1-10-unknown,g10-t2-legacy-cell-1,g10-t2-legacy-cell-2,g10-t2-legacy-cell-3,g10-t2-legacy-cell-4,g10-t2-legacy-cell-5,g10-t2-legacy-cell-6},
  renderer={G10-T2-C1}
}

\HighlightDeclareReference{G10-T2-introduction}{
  status={available},
  members={g10-t2-intro-outcome,g10-t2-intro-event,g10-t2-intro-variable,g10-t2-intro-support,g10-t2-intro-second,g10-t2-intro-third,g10-t2-intro-transform,g10-t2-intro-cutoff,g10-t2-intro-unknown,g10-t2-intro-aux,g10-t2-intro-wait,g10-t2-intro-ops,g10-t2-intro-mass,g10-t2-intro-cum,g10-t2-intro-p,g10-t2-intro-q,g10-t2-intro-break,g10-t2-intro-rate,g10-t2-intro-lambda,g10-t2-intro-expected,g10-t2-intro-meanpop,g10-t2-intro-secondmoment,g10-t2-intro-variance,g10-t2-intro-sd,g10-t2-intro-size,g10-t2-intro-assume,g10-t2-intro-target}
}

\HighlightDeclareReference{G10-T2-full}{
  status={available},
  members={g10-t2-c1-coordinated-10-outcome,g10-t2-c1-coordinated-10-event,g10-t2-c1-coordinated-10-variable,g10-t2-c1-coordinated-10-support,g10-t2-c1-coordinated-10-second,g10-t2-c1-coordinated-10-third,g10-t2-c1-coordinated-10-mass,g10-t2-c1-coordinated-10-unknown,g10-t2-c1-coordinated-10-ops,g10-t2-c1-coordinated-10-assume,g10-t2-c2-coordinated-10-outcome,g10-t2-c2-coordinated-10-event,g10-t2-c2-coordinated-10-variable,g10-t2-c2-coordinated-10-support,g10-t2-c2-coordinated-10-mass,g10-t2-c2-coordinated-10-cum,g10-t2-c2-coordinated-10-unknown,g10-t2-c2-coordinated-10-cutoff,g10-t2-c2-coordinated-10-ops,g10-t2-c2-coordinated-10-assume,g10-t2-c2-coordinated-10-aux,g10-t2-c3-coordinated-10-event,g10-t2-c3-coordinated-10-variable,g10-t2-c3-coordinated-10-support,g10-t2-c3-coordinated-10-p,g10-t2-c3-coordinated-10-q,g10-t2-c3-coordinated-10-second,g10-t2-c3-coordinated-10-mass,g10-t2-c3-coordinated-10-cum,g10-t2-c3-coordinated-10-expected,g10-t2-c3-coordinated-10-break,g10-t2-c3-coordinated-10-transform,g10-t2-c4-coordinated-10-variable,g10-t2-c4-coordinated-10-support,g10-t2-c4-coordinated-10-second,g10-t2-c4-coordinated-10-p,g10-t2-c4-coordinated-10-q,g10-t2-c4-coordinated-10-size,g10-t2-c4-coordinated-10-mass,g10-t2-c4-coordinated-10-cum,g10-t2-c4-coordinated-10-cutoff,g10-t2-c4-coordinated-10-ops,g10-t2-c4-coordinated-10-assume,g10-t2-c5-coordinated-10-wait,g10-t2-c5-coordinated-10-support,g10-t2-c5-coordinated-10-p,g10-t2-c5-coordinated-10-q,g10-t2-c5-coordinated-10-second,g10-t2-c5-coordinated-10-mass,g10-t2-c5-coordinated-10-cum,g10-t2-c5-coordinated-10-cutoff,g10-t2-c5-coordinated-10-event,g10-t2-c5-coordinated-10-assume,g10-t2-c6-coordinated-10-wait,g10-t2-c6-coordinated-10-support,g10-t2-c6-coordinated-10-p,g10-t2-c6-coordinated-10-q,g10-t2-c6-coordinated-10-target,g10-t2-c6-coordinated-10-second,g10-t2-c6-coordinated-10-mass,g10-t2-c6-coordinated-10-cum,g10-t2-c6-coordinated-10-size,g10-t2-c6-coordinated-10-cutoff,g10-t2-c6-coordinated-10-ops,g10-t2-c6-coordinated-10-assume,g10-t2-c7-coordinated-10-variable,g10-t2-c7-coordinated-10-support,g10-t2-c7-coordinated-10-mass,g10-t2-c7-coordinated-10-cum,g10-t2-c7-coordinated-10-cutoff,g10-t2-c7-coordinated-10-rate,g10-t2-c7-coordinated-10-lambda,g10-t2-c7-coordinated-10-second,g10-t2-c7-coordinated-10-ops,g10-t2-c7-coordinated-10-assume,g10-t2-c8-coordinated-10-variable,g10-t2-c8-coordinated-10-support,g10-t2-c8-coordinated-10-second,g10-t2-c8-coordinated-10-mass,g10-t2-c8-coordinated-10-expected,g10-t2-c8-coordinated-10-meanpop,g10-t2-c8-coordinated-10-secondmoment,g10-t2-c8-coordinated-10-variance,g10-t2-c8-coordinated-10-sd,g10-t2-c8-coordinated-10-transform,g10-t2-c8-coordinated-10-ops,g10-t2-c9-coordinated-10-event,g10-t2-c9-coordinated-10-variable,g10-t2-c9-coordinated-10-support,g10-t2-c9-coordinated-10-wait,g10-t2-c9-coordinated-10-second,g10-t2-c9-coordinated-10-p,g10-t2-c9-coordinated-10-q,g10-t2-c9-coordinated-10-size,g10-t2-c9-coordinated-10-target,g10-t2-c9-coordinated-10-mass,g10-t2-c9-coordinated-10-cum,g10-t2-c9-coordinated-10-rate,g10-t2-c9-coordinated-10-lambda,g10-t2-c9-coordinated-10-ops,g10-t2-c9-coordinated-10-assume,g10-t2-intro-outcome,g10-t2-intro-event,g10-t2-intro-variable,g10-t2-intro-support,g10-t2-intro-second,g10-t2-intro-third,g10-t2-intro-transform,g10-t2-intro-cutoff,g10-t2-intro-unknown,g10-t2-intro-aux,g10-t2-intro-wait,g10-t2-intro-ops,g10-t2-intro-mass,g10-t2-intro-cum,g10-t2-intro-p,g10-t2-intro-q,g10-t2-intro-break,g10-t2-intro-rate,g10-t2-intro-lambda,g10-t2-intro-expected,g10-t2-intro-meanpop,g10-t2-intro-secondmoment,g10-t2-intro-variance,g10-t2-intro-sd,g10-t2-intro-size,g10-t2-intro-assume,g10-t2-intro-target}
}

\HighlightDeclareReference{G10-T3-full}{
  status={pending},
  members={}
}

\HighlightDeclareReference{G10-T2-C10}{
  status={pending},
  members={}
}
